\documentclass[a4paper,fleqn]{cas-dc}
\usepackage[numbers]{natbib}
\usepackage{graphicx}
\usepackage{amsmath,amssymb,amsthm}
\usepackage{booktabs}
\usepackage{longtable}
\usepackage{multirow}
\usepackage{array}
\usepackage{tabularx}
\usepackage{xcolor}
\usepackage{algorithm}
\usepackage{algpseudocode}
\usepackage{microtype}
\usepackage{enumitem}
\usepackage{placeins}
\usepackage{tikz}
\setlength{\dblfloatsep}{16pt plus 2pt minus 2pt}
\setlength{\dbltextfloatsep}{16pt plus 2pt minus 2pt}
\usetikzlibrary{arrows.meta,positioning,fit,calc,shapes.geometric,backgrounds}

\setcounter{secnumdepth}{2}
\emergencystretch=2em
\newcolumntype{Y}{>{\raggedright\arraybackslash}X}

\newtheorem{theorem}{Theorem}
\newtheorem{lemma}{Lemma}
\newtheorem{proposition}{Proposition}
\newtheorem{corollary}{Corollary}
\newtheorem{assumption}{Assumption}
\newtheorem{remark}{Remark}
\newtheorem{definition}{Definition}

\DeclareMathOperator*{\argmin}{arg\,min}
\DeclareMathOperator*{\argmax}{arg\,max}
\newcommand{\PP}{\mathbb{P}}
\newcommand{\RR}{\mathbb{R}}
\newcommand{\calF}{\mathcal{F}}
\newcommand{\calP}{\mathcal{P}}
\newcommand{\calR}{\mathcal{R}}
\newcommand{\calG}{\mathcal{G}}
\newcommand{\Valid}{\mathsf{Valid}}
\newcommand{\chain}{\mathrm{chain}}
\newcommand{\CTHR}{\textsc{CTHR}}
\newcommand{\covyes}{\(\checkmark\)}
\newcommand{\covno}{--}
\newcommand{\ie}{\emph{i.e.}}
\newcommand{\eg}{\emph{e.g.}}
\newcommand{\cf}{\emph{cf.}}
\newcommand{\etal}{\emph{et al.}}
\hypersetup{hypertexnames=false}
\pdfstringdefDisableCommands{\def\\{ }}
\definecolor{cthrBlue}{HTML}{1F4E79}
\definecolor{cthrTeal}{HTML}{2A7F62}
\definecolor{cthrOrange}{HTML}{B45F06}
\definecolor{cthrRed}{HTML}{9C2F2F}
\definecolor{cthrGray}{HTML}{F3F5F7}

\begin{document}
\let\WriteBookmarks\relax
\def\floatpagepagefraction{1}
\def\textpagefraction{.001}

\shorttitle{CTHR for Auditable KG-Driven Optimization}
\shortauthors{Zhu et al.}

\title [mode = title]{CTHR: Constraint-Traceable Hierarchical\\Rule Compilation for Auditable\\KG-Driven Optimization}

\author[1]{Longtao Zhu}
\ead{zhulongtao@stu.scu.edu.cn}

\author[2]{Jun Li}

\author[3]{Jiahong Xie}

\author[5]{Ningwei Bai}
\ead{ningwei.bai.23@ucl.ac.uk}

\author[4]{Hai Jiang}
\ead{jianghai@stu.scu.edu.cn}

\author[4]{Yulong Ji}[orcid=0000-0003-3116-5759]
\cormark[1]
\ead{jyl@scu.edu.cn}

\affiliation[1]{organization={College of Computer Science, Sichuan University},
            city={Chengdu},
            country={China}}
\affiliation[2]{organization={System of Systems and Artificial Intelligence Key Laboratory of SICHUAN Province},
            city={Chengdu},
            country={China}}
\affiliation[3]{organization={AVIC Chengdu Aircraft Design and Research Institute},
            city={Chengdu},
            country={China}}
\affiliation[4]{organization={School of Aeronautics and Astronautics, Sichuan University},
            city={Chengdu},
            country={China}}
\affiliation[5]{organization={Department of Computer Science, University College London (UCL)},
            city={London},
            country={United Kingdom}}
\cortext[cor1]{Corresponding author}

%=============================================================================
\begin{abstract}
In safety-critical domains like flight procedure and building design, knowledge-based 
decision support systems must return decisions that are numerically feasible, consistent 
with source rules, and traceable to evidence. Although design specifications can be 
organized as knowledge graphs (KGs), KGs do not directly provide optimizer-ready feasible 
regions with provenance. To address this, we propose \CTHR{} (Constraint-Traceable 
Hierarchical Rule Compilation), a provenance-preserving symbolic modeling framework 
that transforms KG rules into auditable optimization feasible regions. For a given scenario, 
\CTHR{} resolves six rule relations---applicability, dependency, mutual exclusion, 
override, precedence, and parameter propagation---to construct valid rule structures 
and compile them into exact feasible regions annotated with rule- and document-level 
provenance. These regions are consumed by downstream constrained optimizers, yielding 
decisions with certificates tracing back to KG evidence. Theoretical analysis shows 
\CTHR{} preserves the constraint scope induced by source-rule semantics without 
admitting prohibited decisions or excluding valid compliance paths. Experiments on 
200 tasks from flight procedure and building design show that flat-rule modeling 
often produces source-semantically invalid decisions despite high rule recall. On 
the aviation benchmark, the flat baseline returns invalid decisions for all tasks, 
whereas \CTHR{}-style SCIP reduces invalid cases from 150 to 17 and reaches 88.7\% 
Sem-CSR. On the architecture benchmark, \CTHR{}-style backends reduce invalid cases 
from 11 to 2, reaching 96.0\% Sem-CSR. Solver-backend experiments show that HiGHS, 
CP-SAT/OR-Tools, and SCIP successfully optimize over \CTHR{} feasible cells while 
preserving audit certificates. These results validate \CTHR{} as a provenance-preserving 
KG-to-rule-to-feasible-cell modeling layer rather than a new optimizer. Code and data are available at 
\url{https://github.com/Zhulongtao6/CTHR-Constraint-Traceable-Hierarchical-Rule-Compilation-for-Auditable-KG-Driven-Optimization}. 
\end{abstract}

\begin{highlights}
\item CTHR compiles source-KG rules into auditable feasible cells with rule-level and document-level provenance.

\item CTHR resolves applicability, dependency, exclusion, override, precedence, and parameter propagation before optimization.

\item A 200-task aviation/building benchmark evaluates semantic validity, solver reuse, and audit traceability.
\end{highlights}

\begin{keywords}
Knowledge-based systems \sep Knowledge graphs \sep Defeasible reasoning \sep Constraint compilation \sep Auditable optimization
\end{keywords}

\hfuzz=1000pt
\maketitle
\hfuzz=0.1pt

%=============================================================================
\section{Introduction}
\label{sec:intro}

Knowledge-based decision systems are increasingly being deployed in
high-assurance industrial processes, such as flight procedure design and
building design. In these domains, a common practice is to organize
requirements from regulatory documents, standards, and operating manuals as
knowledge graphs (KGs) \cite{horrocks2004swrl,pan2017exploiting}, where
entities, relations, source-document links, and rule annotations jointly
represent domain knowledge. Such KGs support rule retrieval, association among
related requirements, and provenance tracking. However, they do not by
themselves transform the rules retrieved for a concrete decision problem into
the feasible region required for optimization, nor do they directly support
tracing an optimized decision back to the source rules that justify it.

Once domain knowledge has been encoded as a KG, it still needs to be converted
into an operational decision-making mechanism. In practical settings, this
conversion involves mapping KG evidence into a structured rule library,
selecting the rules applicable to a concrete scenario, compiling the resulting
constraints, and solving an optimization problem. At present, this process is
largely carried out by experts who understand both domain regulations and
optimization modeling. Moreover, real-world KGs do not consist of merely flat
rules. They often contain exception rules, cross-document precedence,
mutually exclusive alternatives, and scenario-dependent conditions that can
simultaneously affect multiple requirements. Therefore, transforming symbolic
knowledge into a feasible decision space is a core knowledge-engineering
problem rather than a simple preprocessing step.

Most constrained learning and optimization methods assume a flat
rule-to-constraint mapping: each active rule contributes an independent
inequality, and the conjunction of these inequalities is treated as the
feasible set. Source KGs for technical specifications do not operate in this
way. Exception rules may override base rules for specific procedure categories;
cross-document precedence determines which source should govern when documents
conflict; and alternative templates may correspond to distinct compliance
paths rather than requirements that must be merged. For example, in
building-code compliance, whether a building is equipped with an automatic
sprinkler system may change maximum travel-distance limits, exit-width
thresholds, and fire-separation requirements. If the modeling layer ignores
these rule interactions, an optimizer may solve a region that appears
numerically feasible but is inconsistent with the source-rule semantics. It may
also lose the provenance required to explain why the returned decision is
compliant.

This paper focuses on the modeling layer between a source KG and an optimizer.
Pure neural constrained-optimization approaches \cite{donti2021dc3,navon2021learning}
can improve objective performance under learned or given constraints, but they
usually assume that the feasible set is already correct and do not preserve
rule-level provenance. Pure symbolic reasoning and SMT-based checking methods
\cite{de2008z3} can verify compliance, but they do not directly provide a
source-linked feasible region that can be reused by downstream optimizers
across decision queries. Knowledge-based decision systems therefore require a
modeling method that can transform source-KG rule semantics into optimization
constraints for a concrete problem while preserving the feasible decision scope
and retaining audit evidence.

To address this question, this paper proposes \CTHR{} (Constraint-Traceable Hierarchical
Rule Compilation), a provenance-preserving KG-to-rule-to-feasible-cell modeling
layer for KG-driven optimization. In \CTHR{}, source-KG evidence is first mapped
into a rule library. Then, for each concrete scenario, \CTHR{} resolves rule
applicability, dependencies, mutual exclusions, overrides, and precedence
relations, together with parameter propagation; compiles the remaining valid
rule structures into auditable feasible cells; and preserves source-document
provenance for each feasible cell. The
six rule relations are not introduced as new logical primitives. The core
contribution is their joint use as a source-linked compilation target: \CTHR{}
turns scenario-specific KG rules into solver-consumable feasible cells without
losing alternative compliance paths or provenance.
These feasible cells can then be passed to an exact solver, a classical
constrained optimizer, or a multi-objective search method. \CTHR{} is therefore not proposed
as a new optimizer; its role is to define the source-consistent feasible region
and the traceability interface that downstream solvers operate on. The system
returns not only a numerical decision, but also a source-linked audit record
that allows domain experts to trace the decision back to the active rules and
source documents that support it.

The main contributions of this paper are summarized as follows:

\begin{itemize}[leftmargin=*]
\item \textbf{Joint compilation of KG rule interactions.}
This paper defines how scenario applicability, dependencies, mutual exclusions,
overrides, precedence, and parameter propagation are jointly grounded for a
concrete design scenario. \CTHR{} resolves these source-linked relations before
optimization and represents alternative compliance paths as distinct valid rule
structures rather than flattening them into one constraint set.

\item \textbf{Source-semantics-preserving feasible-cell compilation.}
The framework compiles each valid rule structure into an auditable feasible cell and proves
that the union of feasible cells preserves the feasible region induced by the
source-rule semantics. Each cell retains rule-level and document-level
provenance, so the feasible region is both executable and source traceable.

\item \textbf{Solver-agnostic optimization interface.}
\CTHR{} defines the auditable feasible cells and exposes them to downstream
backends including the default optimizer, ASP/clingo, CP-SAT/OR-Tools, and
SCIP. The solver backend determines optimization strength, while \CTHR{}
provides the source-consistent modeling layer and active-cell interface.

\item \textbf{Audit certificate and provenance traceability.}
For every returned decision, \CTHR{} records the satisfied cell, instantiated
constraints, active valid rule chain, and source provenance. The certificate is
generated from the feasible cell and rule chain, not from a solver-specific
proof log or a post-hoc natural-language explanation.

\item \textbf{Two-domain 200-task benchmark.}
This paper constructs and evaluates a systematic benchmark covering 150 aviation
flight-procedure design tasks and 50 architecture/building-code design tasks.
The benchmark separates method-visible inputs from hidden evaluation labels and
supports evaluation of valid-rule identification, source-semantic decision
validity, solver-backend reuse, rule-library quality, and audit traceability.
\end{itemize}

%=============================================================================
\section{Related Work}
\label{sec:related}

\subsection{Knowledge-Based Decision Support and Provenance}
Knowledge-based decision support in regulated domains has long treated rules as executable objects rather than only retrievable text. Logic-program encodings of legislation, automated rule-based building-design checking, and provenance standards show that compliance systems must preserve both source meaning and audit evidence \cite{sergot1986british,eastman2009automatic,moreau2013provdm}. Recent KBS and neurosymbolic systems further combine symbolic structure with learned components for source-grounded reasoning \cite{amador2024neurosymbolic,some2025comprehensive}. LLM/KG and KG/RAG studies show the value of explicit relational knowledge for organization, retrieval, graph completion, and explanation \cite{pan2024llmkgroadmap,dai2025llmkg,yang2024cpkgc,jin2024llmgraphs,huang2025promptkg}. Similar source-grounded pipelines have appeared in legal, automated-driving, construction, manufacturing, and aviation workflows \cite{hindi2025legalrag,hussien2025roadusers,lee2024constructionsafety,wan2025hragmanufacturing,xiong2025drrag}. These systems strengthen source-grounded access and explanation, but they usually retrieve, validate, rank, or explain knowledge rather than compile provenance-rich rules into auditable optimization feasible regions. \CTHR{} targets this missing modeling layer.

\subsection{KG-to-Rule and LLM-Assisted Rule Modeling}
KGs and declarative rule or constraint languages such as OWL/SWRL and SHACL provide natural representations for structured requirements, graph validation, and rule-like integrity constraints \cite{horrocks2004swrl,pan2017exploiting,knublauch2017shacl}. LLM-assisted KG processing further reduces the knowledge acquisition burden by extracting candidate rules, relations, and provenance fields from technical corpora, as shown in KG-enhanced prompting, knowledge-based answer generation, hybrid RAG, and domain KG construction \cite{liu2025adaptqa,zhou2026ichrag,chen2026dairygoatqa}. However, graph validation and rule extraction do not by themselves produce an optimization model: applicability, dependency, exclusion, override, precedence, and scenario parameters still determine which rules form a valid feasible region. \CTHR{} therefore treats LLM/KG extraction as the front end and performs deterministic symbolic compilation before optimization.

\subsection{Defeasible Reasoning and Symbolic Verification}
Defeasible and non-monotonic reasoning are especially relevant to regulatory domains, where exceptions, priorities, and contextual specificity determine which rules survive \cite{prakken2004three,bonatti2015defeasible}. ASP and SMT solvers provide rigorous machinery for defaults, consistency checking, and formal verification \cite{brewka2011answer,de2008z3}. Neurosymbolic and KBS methods also use rules or graph structure as regularizers, differentiable logic, or relational representations for learning systems \cite{giunchiglia2022road,yang2020neurasp,marra2025relational,chen2024deduction}. These methods are powerful reasoning or verification tools, but they do not by themselves specify how source-linked KG evidence should be grounded, filtered, resolved, and translated into reusable continuous feasible cells. \CTHR{} fills this role as a provenance-preserving compiler rather than a soft rule regularizer or representation learner.

\subsection{Constrained and Multi-Objective Optimization}
Constrained neural optimization and differentiable optimization layers improve feasibility or projection when a constraint set is already given \cite{donti2021dc3,amos2017optnet,agrawal2019differentiable,tordesillas2023rayen}. Preference-conditioned and evolutionary methods search for Pareto-efficient decisions under specified trade-offs \cite{navon2021learning,mahapatra2020multi,deb2002nsga2,li2019two}, and recent LLM-assisted planning or optimization work similarly combines generated procedural knowledge with formal execution or constrained search \cite{liu2026cart,yang2026lmcoea}. These methods are downstream optimizers: they assume that the feasible region is already semantically correct. \CTHR{} addresses the upstream KBS problem by compiling KG-induced rule semantics into the feasible region before any solver optimizes over it.

Taken together, prior work covers important fragments of the problem. Defeasible logics formalize priority and override; policy and compliance languages reason about permissions, prohibitions, obligations, and violations; planning, configuration, and constraint-generation methods instantiate rules into executable models; and ASP/SMT encodings can provide semantic guarantees for their target programs. \CTHR{} targets a different compilation boundary. It starts from source-linked KG rule records whose active subset is scenario-dependent, preserves mutually exclusive and alternative compliance paths as separate valid rule structures, emits optimization-ready feasible cells over design variables, and attaches a certificate linking each optimized decision to rule and document provenance.

%=============================================================================
\section{Motivation}
\label{sec:motivation}

In many high-assurance industrial settings, engineers do not merely face the
problem of whether domain knowledge is available; they face the problem of how
that knowledge can be operationalized. Flight procedure design and building
code design provide representative examples. Domain specifications can be
organized as KGs, where rules, clauses, conditions, and source links can be
stored and retrieved. However, when an engineer needs to solve a concrete
design problem, the KG usually only indicates which rules may be relevant. It
does not directly tell an optimizer which variables may take which values,
which rules are active in the current scenario, which rules should be excluded
because of exceptions, mutual exclusions, or precedence, or what the resulting
feasible region should be. In current practice, this transformation is often
performed manually by a small number of experts who understand both the domain
specification and optimization modeling. The motivation of this paper is to
reduce this gap: KG rules should not remain only queryable knowledge, but
should be transformable into executable, auditable, and optimizable feasible
regions for concrete design scenarios.

The rapid progress of LLMs creates new opportunities for industrial knowledge
modeling. LLMs can help read technical specifications, interpret domain
terminology, summarize rule structure, and integrate knowledge across sources,
making it more realistic to extract structured rule libraries from KG evidence.
However, in safety-critical domains such as aviation and building design, final
decisions should not be produced directly by probabilistic generation. They
should be produced by deterministic, verifiable, and auditable mechanisms. A
more appropriate division of labor is to use LLMs for the tasks at which they
are effective, such as understanding text, organizing knowledge, and generating
candidate rules, while assigning final rule resolution, constraint compilation,
feasibility validation, and certificate generation to a formal symbolic
framework. This view is consistent with LLM-Modulo planning
\cite{kambhampati2024llmmodulo}, which argues that LLMs can help planning by
providing knowledge and candidate structures, but planning, verification, and
execution should remain grounded in external formal modules. \CTHR{} follows
the same principle for KG-driven optimization: LLMs and KGs assist rule-library
construction, while a traceable symbolic rule compiler guarantees the
verifiability and auditability of the final decision.

%=============================================================================
\section{Problem Formulation}
\label{sec:problem}

To formalize the semantic gap between symbolic knowledge and continuous geometric boundaries, a structured regulatory rule base is extracted from the source KG. This computational rule base encapsulates source-linked rule records, document-level provenance, and six explicit non-monotonic rule-interaction mechanisms: applicability, dependency, mutual exclusion, exception override, cross-source precedence, and scenario-dependent parameter propagation. Structurally, this rule base serves as the rigid deterministic input for subsequent defeasible reasoning and polyhedral constraint compilation.

\begin{definition}[Regulatory rule base]
A \emph{regulatory rule base} $\calR$ induced by a source KG over a multi-dimensional scenario space $\Omega$ and a continuous decision space $\RR^n$ consists of a finite set of source-linked regulatory assertions. Each rule $r\in\calR$ is mapped to a parameterized consequent constraint $c_r(x,\eta(\omega)) \equiv \langle a_r(\eta(\omega)),x\rangle \leq b_r(\eta(\omega))$ and enriched with provenance metadata. Rules may be partitioned into disjoint conflict classes when their respective consequents represent alternative or mutually inconsistent compliance paths. The six core rule-interaction mechanisms are formally operationalized as follows:
\begin{enumerate}[label=(\roman*)]
  \item An \emph{applicability} function $\phi_r:\Omega\to\{0,1\}$ defined for each rule, where $\phi_r(\omega)=1$ dictates that rule $r$ is triggered within the grounded scenario $\omega$;
  \item A directed \emph{dependency} relation $\mathsf{dep}\subseteq\calR\times\calR$, where $(r_1,r_2)\in\mathsf{dep}$ denotes that the activation of rule $r_1$ logically requires the concurrent fulfillment of rule $r_2$;
  \item A symmetric \emph{exclusion} relation $\mathsf{excl}\subseteq\calR\times\calR$, where $(r_1,r_2)\in\mathsf{excl}$ stipulates that $r_1$ and $r_2$ cannot simultaneously coexist in any valid deployment;
  \item An \emph{exception-override} relation $\mathsf{ovr}\subseteq\calR\times\calR$, where $(r_e,r_b)\in\mathsf{ovr}$ indicates that a more specific contextual rule $r_e$ defeats a general base rule $r_b$ when both are triggered;
  \item A scenario-dependent \emph{precedence} strict partial order $\succ_\omega$ over conflicting rules within the same conflict class, induced by cross-document hierarchy or source specificity;
  \item A \emph{parameter propagation function} $\eta:\Omega\to\RR^p$ that maps high-dimensional scenario features to contextual parameter vectors. These parameters are not static empirical constants; instead, they are dynamically instantiated by the grounded scenario $\omega$, thereby determining the precise numerical coefficients and affine boundaries utilized in downstream constraints.
\end{enumerate}
\end{definition}

Here, $\omega\in\Omega$ denotes a grounded domain scenario, and $x\in\RR^n$ denotes a candidate decision vector within the engineering optimization domain. The coefficient vector $a_r(\eta(\omega))\in\RR^n$ and the scalar bound $b_r(\eta(\omega))\in\RR$ define the precise algebraic boundaries contributed by rule $r$ after scenario-context grounding, where $\langle\cdot,\cdot\rangle$ denotes the standard Euclidean inner product. The integer $p$ represents the dimension of the context-dependent parameter vector returned by $\eta$. Relational provenance is thoroughly indexed to preserve audit evidence, though it is excluded from the operational rule-interaction typologies.

\begin{definition}[Defeasible resolution]
For any grounded scenario $\omega\in\Omega$, let $\mathrm{App}(\omega) = \{r\in\calR : \phi_r(\omega)=1\}$ denote the active candidate rule set. Within each designated conflict class, an applicable rule $r_e$ is said to \emph{defeat} $r_b$ if and only if $(r_e,r_b)\in\mathsf{ovr}$ or $r_e \succ_\omega r_b$ holds. The \emph{valid rule set} $\Valid(\omega) \subseteq \mathrm{App}(\omega)$ preserves all non-defeated applicable rules after the deterministic resolution of all override and precedence conflicts. Consequently, each maximal subset of $\Valid(\omega)$ that remains pairwise consistent with $\mathsf{excl}$ and topologically closed under $\mathsf{dep}$ synthesizes a self-contained source-level extension, written as a \emph{valid rule structure} $\chain(k)$. Let $\mathcal{K}_{\mathrm{src}}(\omega) = \{1,\ldots,K(\omega)\}$ be the ordered index set enumerating these alternative source-level extensions, satisfying $\chain(k)\subseteq\Valid(\omega)$ for every $k\in\mathcal{K}_{\mathrm{src}}(\omega)$. This extension set is defined only from the source rule theory $(\calR,\phi,\mathsf{dep},\mathsf{excl},\mathsf{ovr},\succ_\omega,\eta)$ and does not depend on any feasible cell emitted by \CTHR{}.
\end{definition}

In this formal resolution scheme, $r_e$ characterizes the overriding or context-specific rule, $r_b$ represents the defeated base rule, and $\succ_\omega$ tracks the dynamic priority relation mandated by cross-document authority boundaries. The integer $K(\omega)$ denotes the total number of alternative, self-consistent compliance structures yielded under scenario $\omega$.

\begin{definition}[Source rule-theory feasible semantics]
The \emph{source rule-theory feasible region} $F_{\mathrm{src}}(\omega)$ for a designated scenario $\omega$ is defined directly from the source-level extensions of the rule theory:
\begin{equation}
  \begin{aligned}
  F_{\mathrm{src}}(\omega)
  = \bigl\{x\in\RR^n \;\big|\;&
  \exists k\in\mathcal{K}_{\mathrm{src}}(\omega):\\
  & \forall r\in\chain(k),\;
  c_r(x,\eta(\omega)) \text{ holds}\bigr\}.
  \end{aligned}
  \label{eq:Fsrc}
\end{equation}
When emphasizing its KG origin, this independently defined source semantics is also written as $F_{\mathrm{KG}}(\omega)$.
\end{definition}

As formulated in \eqref{eq:Fsrc}, the global optimization space is represented not as a naive flat conjunction of all applicable requirements, but as a disjunctive semantic structure. A candidate decision vector $x$ is deemed semantically valid if it satisfies the complete conjunction of constraints within at least one isolated valid rule structure $\chain(k)$. This formulation ensures that alternative compliance templates or mutually exclusive regulatory pathways are preserved as distinct geometric cells, preventing the spurious over-constraint or invalid optimization paths common in flat compilation baselines. Figure~\ref{fig:flat_vs_cthr_example} illustrates this distinction using a building-code example with alternative roll-in shower templates.

\begin{figure*}[t]
\centering
\includegraphics[width=\textwidth]{"figures/figs/fig. 2. an example.jpg"}
\caption{Illustrative comparison between conventional flat modeling and \CTHR{} valid-structure modeling. Flat modeling may merge mutually exclusive compliance templates into a single constraint set, whereas \CTHR{} preserves alternative valid rule structures and compiles them into a disjunctive semantic feasible region.}
\label{fig:flat_vs_cthr_example}
\end{figure*}
\FloatBarrier

For a concrete engineering design query, the task provides a grounded scenario, decision variables, and objectives, while \CTHR{} supplies the source-preserving feasible region $F_{\mathrm{KG}}(\omega)$ and its provenance metadata. The downstream optimizer then searches within this emitted region, and the returned decision is evaluated against both the constructed feasible region and the hidden source-rule semantics in Section~\ref{sec:exp}. Thus, the problem formulation isolates the contribution of \CTHR{} to semantic feasible-region construction rather than to a particular objective scalarization or solver backend.

%=============================================================================
\section{CTHR: Methodological Framework}
\label{sec:method}

The CTHR framework establishes a provenance-preserving knowledge-engineering
pipeline that turns source-linked technical rules into problem-specific,
auditable feasible cells for optimization. As shown in Figure~\ref{fig:arch},
the CTHR method comprises four parts: Knowledge Acquisition, CTHR Modeling
Layer, Solver Backends, and Optimized Decision \& Audit Certificate.

The Knowledge Acquisition part constructs the knowledge graph and rule library.
The CTHR Modeling Layer resolves task-valid rules and compiles them into
source-traceable feasible cells. Solver backends then optimize the task
objective within these cells; they are interchangeable execution engines rather
than rule interpreters. Finally, the Optimized Decision \& Audit Certificate
component links each returned decision back to the active cell, valid rule
chain, and source documents.

\begin{figure*}[t]
\centering
\includegraphics[width=\textwidth]{"figures/figs/fig. 1. overview.jpg"}
\caption{Overall automated end-to-end CTHR framework topology, mapping the pipeline from raw technical specifications to optimized decisions and audit certificates.}
\label{fig:arch}
\end{figure*}


\subsection{Knowledge Acquisition}
\label{sec:knowledge_acquisition}

Knowledge Acquisition converts unstructured technical specification documents into computable, traceable rule-based knowledge representations. The CTHR framework utilizes Cognee to construct a source-linked knowledge graph, following the broader trend of using LLM-assisted knowledge engineering and KG/RAG pipelines to organize domain documents into structured, reusable knowledge assets \cite{cognee2025,some2025comprehensive,zhou2026ichrag,wan2025hragmanufacturing}. This process parses raw PDFs into granular text chunks while faithfully retaining source provenance metadata, including document names and page numbers, for each extracted fragment. Subsequently, Cognee executes the structured extraction of regulatory entities, applicability criteria, numerical thresholds, measurement units, design objects, constraint relations, and cross-clause references from these text chunks, mapping them into a unified graph topology of nodes and edges. The resulting knowledge graph preserves both the literal normative content and the explicit provenance links connecting regulatory evidence to the original documents, thereby enabling downstream rules to be deterministically traced back to their specific source documents, clauses, or foundational evidence.

Following the construction of the knowledge graph, CTHR employs a Large Language Model (LLM) to transform the knowledge graph into a structured rule library. Operating as a KG-to-rule extractor, the LLM converts graph-structured evidence into uniformly formatted rule records, where each rule record encapsulates the original source text fragment, document identifier, section or page number prompts, associated entity nodes, numerical values, measurement units, and the local graph topology of adjacent relations.

The extraction prompt design for the LLM is structured around four core elements: first, task instructions that restrict the model to extracting technical specifications derived strictly from the provided graph evidence; second, input evidence comprising the source text chunk, related entities, numerical values, units, and local graph adjacency information; third, a target output schema specifying that every rule record must contain a rule identifier, rule name, domain, rule type, applicability guard, consequent constraints, parameters, inter-rule relations, and provenance fields; fourth, formatting constraints mandating a valid, parsable JSON structure while strictly prohibiting the fabrication of thresholds, units, or inter-rule relationships lacking explicit source evidence.

The generated rule records undergo automated verification and normalization to guarantee that their JSON format, rule identifiers, constraint expressions, measurement units, numerical sources, and provenance metadata fully satisfy subsequent symbolic modeling requirements. The validated rules are then compiled into a global rule library:
\begin{equation}
\mathcal{R}=\{r_i\}_{i=1}^{N}.
\end{equation}
Each rule $r_i$ encapsulates two distinct categories of information: the first is executable semantic information, such as applicability conditions, parameters, constraint expressions, and rule typologies; the second is structural and audit attributes, including dependency, exclusion, override, precedence, and parameter propagation relationships, supplemented by both document-level and rule-level provenance data.

\subsection{CTHR Modeling Layer}
\label{sec:cthr_modeling_layer}

\subsubsection{Candidate Rule Grounding}
\label{sec:candidate_grounding}

The objective of candidate rule grounding is to generate a high-recall candidate rule set $C(q)$ from the global rule library $\mathcal{R}$ for a specific optimization query $q$. Specifically, the framework computes a deterministic task-grounding score for each rule $r_i$ in the rule library. This score comprehensively aggregates multiple matching metrics: the lexical token overlap between the optimization query text and the rule attributes, such as rule name, rule identifier, rule type, and the provenance section; the semantic alignment between fields and values within the applicability guards and the grounded scenario facts; the structural mapping of constraint variables to task decision variables; the exact consistency of measurement units; and the alignment of the originating document or source domain with the current task scope. For rules featuring explicit applicability guards, an additional weight is awarded if the guard conditions are satisfied within the current scenario. This scoring procedure is formulated as follows:
\begin{equation}
s(r_i,q) = s_{\mathrm{text}} + s_{\mathrm{guard}} + s_{\mathrm{var}} + s_{\mathrm{unit}} + s_{\mathrm{prov}}.
\end{equation}

Where $s_{\mathrm{text}}$ represents the lexical overlap component, computed by normalizing the visible task fields and rule metadata into distinct token sets and evaluating their intersection size. This normalization pipeline involves lowercasing, tokenization via underscores, hyphens, and punctuation, removal of standard stopwords, and synonym expansion for domain-specific abbreviations. The remaining score components are derived from guard satisfaction, variable binding, unit consistency, and source document matching, respectively. Consequently, the task-grounding score operates as a deterministic, logic-driven symbolic matching metric rather than an embedding-based probabilistic vector similarity.

Subsequently, the framework ranks the rules based on their aggregated task-grounding scores and retains high-scoring records as initial candidates. Crucially, the output of this candidate grounding phase does not represent the final valid rule set, but rather an intentionally broad candidate set:
\begin{equation}
C(q) = \{r_i \in \mathcal{R} \mid s(r_i,q) \ge \tau \ \text{or}\ r_i\ \text{passes recall guard}\}.
\end{equation}
This candidate set is deliberately over-inclusive and is typically significantly larger than the final resolved valid rule set.

\subsubsection{Rule-Interaction Resolution}
\label{sec:rule_interaction_resolution}

Candidate rule grounding favors recall, so $C(q)$ usually contains both active rules and rules recalled only through superficial matches in vocabulary, variables, or provenance. Some recalled rules also affect the final rule topology through dependency, exception, exclusion, precedence, or parameter propagation rather than through a direct numerical constraint. Rule-interaction resolution therefore contracts the broad candidate set into the valid rule set for the current task:
\begin{equation}
V(q) \subseteq C(q).
\end{equation}
Algorithm~\ref{alg:rule_interaction_resolution} summarizes the procedure. It first applies task-profile filtering, then optionally invokes an LLM-assisted filter over public task and rule fields, and finally resolves explicit rule relations symbolically.

\begin{algorithm}[t]
\caption{Rule-Interaction Resolution}
\label{alg:rule_interaction_resolution}
\begin{algorithmic}[1]
\Require Task $q$; broad candidate set $C(q)$; rule relation graph $\mathcal{E}_{\mathrm{rel}}$; optional LLM filter $\mathcal{M}$
\Ensure Valid rule set $V(q)$

\State Construct task profile $P(q) = (G_q^{\mathrm{req}}, G_q^{\mathrm{allow}}, G_q^{\mathrm{block}}, \sigma_q)$
\State $C_{\mathrm{seed}}(q) \leftarrow \emptyset$

\For{each $r \in C(q)$}
    \State Map rule $r$ to rule groups $G(r)$
    \State Compute profile indicators $I_{\mathrm{req}}, I_{\mathrm{allow}}, I_{\mathrm{block}}, I_{\mathrm{unmatch}}$
    \State Compute visible binding score $v(r,q)$ from guard fields, decision variables, units, and normalized tokens
    \State Compute profile score: $s_{\mathrm{prof}}(r,q) = w_{\mathrm{req}}I_{\mathrm{req}} + w_{\mathrm{allow}}I_{\mathrm{allow}} - w_{\mathrm{block}}I_{\mathrm{block}} - w_{\mathrm{unmatch}}I_{\mathrm{unmatch}} + v(r,q)$
    \If{$s_{\mathrm{prof}}(r,q) \ge \tau_{\mathrm{seed}}$}
        \State $C_{\mathrm{seed}}(q) \leftarrow C_{\mathrm{seed}}(q) \cup \{r\}$
    \EndIf
\EndFor

\State Remove blocked or unmatched rules from $C_{\mathrm{seed}}(q)$
\State Re-inject required rules that are not blocked
\State Obtain $C_{\mathrm{prof}}(q)$

\If{LLM-assisted filtering is enabled}
    \State $C_{\mathrm{prof}}(q) \leftarrow \mathcal{M}(q, C_{\mathrm{prof}}(q))$
\EndIf

\State $V(q) \leftarrow C_{\mathrm{prof}}(q)$
\State Remove rules whose applicability guards are false under $q$
\State Close $V(q)$ under dependency and parameter/formula propagation
\State Resolve exclusion, override, and precedence conflicts in $V(q)$

\State \textbf{return} $V(q)$
\end{algorithmic}
\end{algorithm}


For task-profile filtering, CTHR builds
\begin{equation}
P(q)=\left(G_q^{\mathrm{req}},G_q^{\mathrm{allow}},G_q^{\mathrm{block}},\sigma_q\right),
\end{equation}
where $G_q^{\mathrm{req}}$, $G_q^{\mathrm{allow}}$, and $G_q^{\mathrm{block}}$ are rule-group tags required, allowed, or blocked by the task, and $\sigma_q$ enables strict filtering. Each candidate rule is mapped to group tags $G(r)$ using its identifier, name, type, source domain, guards, variables, and provenance. This mapping induces four operational states: \emph{required} rules are explicitly triggered by the task; \emph{allowed} rules remain compatible with the task domain or structure; \emph{blocked} rules conflict with the scenario, design object, or domain; and \emph{unmatched} rules have neither profile support nor structural support.

As shown in Algorithm~\ref{alg:rule_interaction_resolution}, these states are converted into a profile score. Required and allowed matches receive positive weights, blocked and unmatched states receive penalties, and the visible binding term $v(r,q)$ captures guard-field overlap, bindable decision variables, unit consistency, and normalized token overlap. In the main experiments, the strict-profile configuration fixes $w_{\mathrm{req}}=7$, $w_{\mathrm{allow}}=3$, $w_{\mathrm{block}}=9$, $w_{\mathrm{unmatch}}=5$, and $\tau_{\mathrm{seed}}=8.0$; the last two weights enter the score with the minus signs shown in Algorithm~\ref{alg:rule_interaction_resolution}. Rules that pass the seed threshold form $C_{\mathrm{prof}}(q)$ after blocked/unmatched removal and required-rule re-injection. The optional LLM-assisted filter only uses public task descriptions, scenario facts, decision variables, candidate rule text, guards, and provenance fields; reference valid rules and solver outcomes are never exposed.

The final step performs symbolic resolution over $C_{\mathrm{prof}}(q)$. CTHR removes rules whose guards are false, closes the set under dependency and parameter/formula propagation, and resolves exclusion, override, and precedence conflicts. The output is
\begin{equation}
V(q)=\mathrm{Resolve}\left(C_{\mathrm{prof}}(q),\mathcal{E}_{\mathrm{rel}},\phi_q\right),
\end{equation}
where $\mathcal{E}_{\mathrm{rel}}$ is the rule-relation graph and $\phi_q$ contains the task scenario facts, bindings, and unit profile.


\subsubsection{Valid Rule Structure Construction}
\label{sec:valid_rule_structure_construction}

Rule-interaction resolution returns a valid rule set $V(q)$, but this set is still flat. Technical specifications may contain exceptions, parameter definitions, formulas, and alternative compliance paths that should not be merged into one conjunction. CTHR therefore organizes $V(q)$ into self-consistent valid rule structures:
\begin{equation}
\mathcal{S}(q) = \{S_1, S_2, \ldots, S_{K(q)}\},
\end{equation}
where each $S_k$ is independently compiled into one feasible cell.

CTHR first constructs the induced relation graph
\begin{equation}
H_q = (V(q), \mathcal{E}_{\mathrm{rel}}|_{V(q)}),
\end{equation}
where $\mathcal{E}_{\mathrm{rel}}|_{V(q)}$ preserves the dependencies, exclusions, overrides, precedence orders, and parameter-propagation links among valid rules. Rules are then separated into core rules that must be retained and branch rules governed by alternative, exclusion, override, or precedence relations. Conflict classes are denoted as
\begin{equation}
\mathcal{B}(q) = \{B_1, B_2, \ldots, B_m\},
\end{equation}
where each $B_j$ contains rules or rule combinations that cannot co-exist. Override and precedence remove defeated branches; unresolved but mutually valid alternatives are retained as separate structures.

For each candidate structure, CTHR enforces dependency closure and parameter propagation:
\begin{equation}
\forall r \in S_k, \quad \mathrm{Dep}(r) \cap V(q) \subseteq S_k.
\end{equation}
Scenario-dependent parameters are instantiated from the grounded task facts. Structures with unsatisfiable guards, missing dependencies, or mutually exclusive rules are discarded.

The resulting valid rule structures are formally represented as:
\begin{equation}
S_k = (R_k, E_k, \gamma_k, \Pi_k),
\end{equation}
where $R_k$ is the rule subset, $E_k$ is the induced relation subgraph, $\gamma_k$ is the activation condition, and $\Pi_k$ stores rule- and document-level provenance. Each $S_k$ must be dependency-closed and exclusion-consistent. Downstream compilation then processes each $S_k$ independently, preserving the disjunctive semantics of alternative regulatory pathways.

\subsubsection{Feasible-Region Compilation}
\label{sec:feasible_region_compilation}

Valid rule structure construction yields $\mathcal{S}(q)=\{S_1,\ldots,S_{K(q)}\}$, but each symbolic structure must still be converted into executable constraints. CTHR compiles every $S_k$ into a feasible cell and aggregates these cells into the task-level search space.

For each rule $r$, the library stores compiled rule-to-constraint templates $\mathcal{T}(r)=\{t_{r,1},t_{r,2},\ldots\}$. A template contains the constraint expression, required symbols, parameters, units, constraint role, and provenance mapping. At runtime, CTHR binds these symbols to the decision variables, bounds, unit system, and scenario parameters of task $q$.

For $S_k=(R_k,E_k,\gamma_k,\Pi_k)$, the compiler verifies required symbols, performs variable/unit/parameter binding, and merges the instantiated constraints with the baseline task domains:
\begin{equation}
P_k(q) = \left\{ x \in \mathcal{X}_q \;\middle|\; \bigwedge_{r \in R_k} \bigwedge_{c \in \Gamma(r,q)} c(x;\eta_q) \text{ holds} \right\},
\end{equation}
where $\mathcal{X}_q$ is the baseline task domain, $\Gamma(r,q)$ is the executable constraint set emitted from rule $r$, and $\eta_q$ contains propagated parameter values.

Alternative structures are kept disjunctive rather than flattened:
\begin{equation}
F_{\mathrm{CTHR}}(q) = \bigcup_{k=1}^{K(q)} P_k(q).
\label{eq:fcthr}
\end{equation}
Thus, a decision is compliant if it satisfies one complete feasible cell, which preserves alternatives and exceptions that flat conjunction would merge incorrectly.

Each cell also stores audit metadata:
\begin{equation}
\mathrm{cell}_k = (\mathrm{cell\_id}, R_k, \Gamma_k, \gamma_k, \Pi_k, \mathcal{B}_k),
\end{equation}
where $\Gamma_k$ contains instantiated constraints and $\mathcal{B}_k$ records bindings and parameter instantiations. These records allow a later solution to be mapped back to the active cell, valid rule structure, constraint templates, and source evidence. Overall, feasible-region compilation performs
\begin{equation}
\mathcal{S}(q) \longrightarrow \{P_1(q), P_2(q), \ldots, P_{K(q)}(q)\} \longrightarrow F_{\mathrm{CTHR}}(q).
\end{equation}
The compiled representation remains solver-agnostic.

\subsection{Solver Backends}
\label{sec:solver_backends}

Following the processing within the CTHR modeling layer, the system outputs a
set of feasible cells rather than a single mathematical programming
model tied to one backend. CTHR is responsible for rule filtering, relational
dependency resolution, valid rule structure construction, and feasible-cell
compilation; the downstream backend is responsible for optimizing within the
emitted search space. For a designated task $q$, the global feasible region
synthesized by CTHR is the feasible-cell union in \eqref{eq:fcthr}, where each
$P_k(q)$ denotes a distinct feasible cell derived from a valid rule structure.
The objective of the downstream solver backend is to search for an optimal decision vector within this search space:
\begin{equation}
\hat{x}_q \in \arg\min_{x \in F_{\mathrm{CTHR}}(q)} g_q(x),
\end{equation}
where $g_q(x)$ denotes the target objective function or a scalarized multi-objective formulation corresponding to the current task profile.

Crucially, the solver backends do not dynamically re-evaluate rule validity,
nor do they re-resolve inter-rule dependencies, mutual exclusions, exception
overrides, or precedence hierarchies. These semantic operations are handled
upstream within the CTHR modeling layer. The downstream solver ingests finalized
feasible cells, variable bounds, structural constraint expressions, and
objective metrics. Consequently, the backends function as optimization engines
rather than rule-interpretation modules.

Nevertheless, different solver backends require varied symbolic representations for the identical set of geometric feasible cells. For the ASP/clingo backend, CTHR translates the rule topology and cell membership rules into logical facts, choice rules, and integrity constraints. For the CP-SAT/OR-Tools backend, constraints are systematically transformed into integer-scaled linear equations, boolean decision variables, and indicator conditions. For the SCIP backend, constraints are mapped into continuous, integer, linear, non-linear, or mixed-integer programming formulations depending on their structural expression types. Thus, while CTHR outputs a unified feasible region semantics, specialized backend adapters are deployed to compile the identical set of cells into the respective target modeling languages. Solver agnosticism is therefore an interface-level claim: CTHR does not automatically make every solver understand the emitted cells. A backend adapter must encode the cell-selection conditions, variables, constraints, bounds, and scalarized objective in the target solver language; solvers whose modeling language cannot express the required constraint class need an additional relaxation, reformulation, or verification layer.

This architectural design makes CTHR solver-agnostic. Default numerical
optimizers, ASP/clingo, CP-SAT/OR-Tools, and SCIP can operate over the same
\CTHR{} feasible cells, while differing in their symbolic encodings, search
mechanisms, and optimization strength. The output synthesized by a backend
contains the optimized decision vector, objective values, and the active cell
that the decision satisfies. CTHR then maps this active cell back to the
corresponding valid rule structure, originating rules, constraint templates,
and document-level source provenance.

In summary, CTHR defines the auditable feasible cells and their provenance
links, while the selected backend determines how strongly the objective is
optimized within those cells. This decoupling lets strong external solvers be
used without requiring them to natively understand KG provenance or defeasible
rule interactions, but each solver family still requires an explicit adapter
for its own modeling language.

\subsection{Optimized Decision and Audit Certificate}
\label{sec:optimized_decision_audit_certificate}

After solver execution, CTHR returns both the optimized decision and an audit certificate. The output is
\begin{equation}
o_q = (\hat{x}_q, g_q(\hat{x}_q), \mathrm{cell}_k, \mathrm{cert}_q),
\end{equation}
where $\hat{x}_q$ is the decision vector, $g_q(\hat{x}_q)$ is its objective value, $\mathrm{cell}_k$ is the satisfied feasible cell, and $\mathrm{cert}_q$ is the certificate.

CTHR first identifies the active cell:
\begin{equation}
\hat{x}_q \in P_k(q).
\end{equation}
Because each cell already stores $R_k$, $\Gamma_k$, $\mathcal{B}_k$, $\gamma_k$, and $\Pi_k$, the solution can be traced through
\begin{equation}
\hat{x}_q \rightarrow P_k(q) \rightarrow S_k \rightarrow R_k \rightarrow \mathrm{Provenance}(R_k).
\end{equation}
The certificate records four fields: the decision summary, the satisfied
instantiated constraints, the active valid rule chain, and the source
provenance. Since provenance links are preserved from document chunks to rule
records and then to constraint templates, certificate generation is
deterministic. It is not a solver proof log and not a post-hoc LLM explanation:
the certificate is assembled from the selected feasible cell and its stored rule
chain. If a constraint $c\in\Gamma_k$ is derived from rule $r$ with provenance
$\Pi(r)=(\mathrm{doc},\mathrm{page},\mathrm{section},\mathrm{chunk},\mathrm{kg\_nodes})$, then the certificate records $c\rightarrow r\rightarrow\Pi(r)$:
\begin{equation}
\mathrm{cert}_q = \left\{ (c, r, \Pi(r)) \;\middle|\; c \in \Gamma_k, \; r \in R_k \right\}.
\end{equation}
This gives domain experts a direct path from a numerical decision to the executable constraints, valid rules, and original technical evidence that justify it.


%=============================================================================
\section{Theoretical Analysis}
\label{sec:theory}

The purpose of this section is to make explicit what must be guaranteed before the compiled knowledge graph (KG) can be safely used for optimization. Empirical results can show that \CTHR{} works on the tested scenarios, but they do not by themselves prove that compilation preserves the constraint scope induced by the source KG. If a compiler retains defeated rules, merges mutually exclusive compliance paths, or ignores precedence, the optimizer may solve a feasible region that is either too large or too small. The theorem below compares \CTHR{} against the independent source rule-theory semantics $F_{\mathrm{src}}(\omega)$ in \eqref{eq:Fsrc}, which is defined before any feasible cell is emitted. This separation avoids using the compiled representation to define the semantic target. To formally bridge this theoretical analysis with the algorithmic framework detailed in Section 5, a grounded scenario $\omega \in \Omega$ is defined as the formal evaluation context mapped directly from the active runtime task query $q$ and its associated scenario facts $\phi_q$.

Theorem~\ref{thm:sp} is the main correctness result for this compilation target. It is not claimed as a general semantic-preservation theorem for all symbolic encodings; rather, it states the property needed by \CTHR{}: the emitted feasible-cell union is sound and complete with respect to the independently defined source rule theory. A proof sketch is given below, and the detailed lemma proofs are provided in Appendix~\ref{app:proofs}.

\begin{assumption}[Well-founded defeat]\label{ass:defeat}
The defeat relation over conflict classes is acyclic and uniquely determined by scenario $\omega$ (no cycles in override/precedence chains after normalization).
\end{assumption}

\begin{assumption}[Acyclic dependencies]\label{ass:dep}
The \textsf{dep} relation is acyclic. Parameter propagation $\eta(\omega)$ is uniquely determined from $\omega$ by topological evaluation.
\end{assumption}

\begin{assumption}[Polytope non-emptiness]\label{ass:nonempty}
Each feasible cell $P_k$ is a nonempty, closed polytope within the bounded design space.
\end{assumption}

In the practical pipeline, the geometric property required by Assumption~\ref{ass:nonempty} is rigorously enforced during the compilation layer. CTHR preemptively invokes the downstream constraint programming backend (specifically, Google OR-Tools CP-SAT) to execute a structural feasibility check, filtering out inherently self-contradictory or unresolvable rule structures before they are emitted to final optimization.

\begin{theorem}[Source-Theory Soundness and Completeness]\label{thm:sp}
Under Assumptions~\ref{ass:defeat}--\ref{ass:nonempty}, let the emitted \CTHR{} feasible region for scenario $\omega$ be
\begin{equation}
  F_{\mathrm{emit}}(\omega)=
  \bigcup_{\substack{k\in\mathcal{K}_{\mathrm{src}}(\omega)\\ \omega\models\gamma_k}} P_k .
  \label{eq:femit}
\end{equation}
For every grounded scenario $\omega\in\Omega$ and decision $x\in\RR^n$, the compiler is:
\begin{align}
\textnormal{sound:}\quad & x\in F_{\mathrm{emit}}(\omega)\Rightarrow x\in F_{\mathrm{src}}(\omega), \label{eq:soundness}\\
\textnormal{complete:}\quad & x\in F_{\mathrm{src}}(\omega)\Rightarrow x\in F_{\mathrm{emit}}(\omega). \label{eq:completeness}
\end{align}
Consequently, $F_{\mathrm{emit}}(\omega)=F_{\mathrm{src}}(\omega)$, and any downstream optimizer constrained to the emitted feasible cells solves over exactly the feasible decision scope induced by the source rule theory.
\end{theorem}

The proof decomposes into five lemmas:

\begin{lemma}[Applicability Correctness, L1]\label{lem:app}
For every scenario $\omega$ and valid rule structure $k$: $\omega\models\gamma_k$ iff every rule in $\mathrm{chain}(k)$ is applicable, the chain is closed under dependencies, and the chain is internally consistent with the exclusion relation. Applicable alternatives outside the chain may still exist and form separate source-level extensions.
\end{lemma}

\begin{lemma}[Override Correctness, L2]\label{lem:override}
Let $r_e$, $r_b$ be applicable rules in the same conflict class with $r_e$ overriding $r_b$. Then for every $k\in\mathcal{K}_{\mathrm{src}}(\omega)$, the valid rule structure for $k$ contains the consequent of $r_e$ and excludes the defeated consequent of $r_b$.
\end{lemma}

\begin{lemma}[Precedence Correctness, L3]\label{lem:prec}
Let $r_i \succ_\omega r_j$ be conflicting applicable rules. Every valid rule structure retains the consequent of $r_i$ and excludes the conflicting consequent of $r_j$. Non-conflicting applicable rules are both retained.
\end{lemma}

\begin{lemma}[Parameter Propagation Correctness, L4]\label{lem:param}
For every propagated rule $r$, scenario $\omega$, and design point $x$: $c_r(x,\eta(\omega)) \iff \hat{c}_{r,\omega}(x)$, where $\hat{c}_{r,\omega}$ is the compiler-instantiated constraint after symbolic substitution of $\eta(\omega)$.
\end{lemma}

\begin{lemma}[Geometry Correctness, L5]\label{lem:geom}
For every feasible cell $k$ with $\omega\models\gamma_k$: $x\in P_k \iff x$ satisfies all surviving instantiated consequents of $\mathrm{chain}(k)$.
\end{lemma}

\begin{proof}[Proof of Theorem~1 (sketch)]
Soundness: Let $x \in F_{\mathrm{emit}}(\omega)$. Then $x\in P_k$ for some $k\in\mathcal{K}_{\mathrm{src}}(\omega)$ with $\omega\models\gamma_k$. By L5, $x$ satisfies all emitted instantiated consequents of $\mathrm{chain}(k)$. By L2--L3, these consequents are exactly the non-defeated source-level consequents retained by the source rule theory. By L4, propagated parameter instantiation is source-equivalent. By L1, $\mathrm{chain}(k)$ is an applicable, dependency-closed, exclusion-consistent source-level extension. Therefore $x$ satisfies the definition of $F_{\mathrm{src}}(\omega)$.

Completeness: Let $x \in F_{\mathrm{src}}(\omega)$. By \eqref{eq:Fsrc}, there exists a source-level extension $\mathrm{chain}(k)$ such that all source-level consequents in that extension hold at $x$. By L1, the compiler activation condition satisfies $\omega\models\gamma_k$. By L2--L3, the compiler removes exactly the defeated alternatives and preserves the same non-defeated extension. By L4, parameter substitution preserves each source consequent. By L5, the instantiated consequents define $P_k$, so $x\in P_k\subseteq F_{\mathrm{emit}}(\omega)$.
\end{proof}

\begin{remark}
Theorem~\ref{thm:sp} is a conditional soundness-and-completeness result relative
to the independently defined source rule-theory semantics in \eqref{eq:Fsrc}.
It does not state that the entire end-to-end pipeline always recovers the
correct source-level extension. Therefore, invalid cases observed in
Section~\ref{sec:overall_perf} should be interpreted as pipeline errors from
candidate-to-valid recovery, rule-library quality, or backend optimization over
the emitted feasible cells, rather than as a contradiction of the theorem
itself.
\end{remark}

%=============================================================================
\section{Experiments}
\label{sec:exp}

The experiments test the central claim of this paper: a source KG must be
compiled through valid rule structures before its rules are used as optimization
constraints. Section~\ref{sec:overall_perf} evaluates end-to-end
KG-to-constraint modeling and solver-backend reuse. Section~\ref{sec:candidate_valid_exp}
then isolates the candidate-to-valid rule recovery step before feasible-region
compilation and optimization. Section~\ref{sec:rule_library_exp} evaluates how
different LLM-generated rule libraries affect the same \CTHR{} modeling
pipeline. Section~\ref{sec:audit_traceability_exp} finally evaluates whether
the optimized decisions returned through different solver backends retain
complete audit certificates. Section~\ref{sec:ablation_exp} studies which
\CTHR{} components are responsible for the observed end-to-end performance, and
Section~\ref{sec:fairness_scalability_exp} reports supplemental fairness and
controlled scalability checks.

\subsection{Experimental Setup}
\label{sec:exp_setup}

\subsubsection*{Domain Corpora and KG Exports}
The experiments use two source-traceable domain corpora. The aviation corpus is derived from the Civil Aviation Administration of China (CAAC) flight-procedure design specification AC-97-FS-005R1. The architecture corpus combines three regulatory sources: the 2010 ADA Standards for Accessible Design, the 2021 International Building Code (IBC), and the 2021 International Fire Code (IFC). Each corpus is converted into a source-linked knowledge graph (KG) using Cognee~\cite{cognee2025}. The exported data fully retains chunk identifiers, entity nodes, KG edges, document labels. Table~\ref{tab:domain_corpora} summarizes the specific statistics of the two KG exports used in the experiments.

\begin{table*}[t]
\centering
\caption{Domain corpora and KG exports used for the benchmark.}
\label{tab:domain_corpora}
\scriptsize
\setlength{\tabcolsep}{3pt}
\begin{tabularx}{\textwidth}{p{0.12\textwidth}Yrr}
\toprule
Domain & Source corpus & KG nodes & KG edges \\
\midrule
Aviation & CAAC flight-procedure design specification AC-97-FS-005R1 (refers to ICAO Doc 8168). & 4,659 & 6,876 \\
Architecture & 2010 ADA Standards for Accessible Design; 2021 International Building Code; 2021 International Fire Code. & 42,350 & 61,239 \\
\bottomrule
\end{tabularx}
\end{table*}

\subsubsection*{Benchmark Tasks}
A controlled source-grounded benchmark is constructed with 200 rule-grounded
design and compliance tasks: 150 aviation flight-procedure design tasks and 50
architecture/building-code design tasks. Task scenarios are curated from
source specifications and domain practice scenarios; they are not generated from
the parsed KG, rule library, feasible cells, or solver results. The aviation
tasks are derived from CAAC AC-97-FS-005R1. To audit aviation reference-label
quality, 100 sampled aviation tasks were reviewed by three flight-procedure
design experts certified through the CAAC Flight Procedure Design Basic
Knowledge Examination; Appendix~\ref{app:aviation_expert_review} reports the
review protocol and scores. The architecture tasks are grounded in the ADA,
IBC, and IFC rule sources and adapt scenario themes from NCARB ARE practice
materials after rewriting, parameterization, and mapping to the architecture
rule library. The architecture benchmark does not reproduce original question
stems, answer options, or official explanations, and was manually reviewed.

Each task stores algorithm-visible inputs, a public scenario model, and
evaluator-only references. Visible inputs include task metadata, scenario facts,
decision variables, objectives or scalarized preferences, and the public
scenario model; hidden references contain surviving rule IDs, executable
rule-derived constraints, and certificate targets. The aviation expert audit
reports aggregate quality scores rather than formal inter-annotator agreement.
Figure~\ref{fig:benchmark_visible_inputs} shows a visible task input from the
two domains.

\begin{figure*}[t]
\centering
\includegraphics[width=\textwidth]{"figures/figs/fig. 3. an example of dataset .jpg"}
\caption{Example algorithm-visible task inputs from the aviation and
architecture benchmarks. The figure shows only the public task metadata,
scenario facts, decision variables, objectives, and public scenario model.
Hidden rule IDs, rule-derived feasible-region constraints, and provenance
targets are used only for evaluation.}
\label{fig:benchmark_visible_inputs}
\end{figure*}


\subsubsection*{Method-Visible Inputs and Hidden Evaluation Labels}
To avoid oracle leakage, every experiment separates method inputs from
evaluation labels. The method-visible inputs include the task domain, task type,
scenario facts, public scenario model, decision variables and units, objectives
or preference weights, the generated rule library, candidate rules available to
the evaluated method, and the rule relations recorded in that library. Hidden
evaluation labels include reference valid rule IDs, reference valid rule
structures, rule-derived semantic constraints, certificate targets, invalid-case
labels, and best-known objectives used for objective-gap reporting. These hidden
fields are read only by the evaluator when computing Rule P/R, Sem-CSR, invalid
cases, certificate completeness, and objective gaps. They are not supplied to
candidate grounding, the profile resolver, the optional LLM filter, feasible
cell compilation, or the solver backends. Experiments that intentionally fix
feasible cells or use oracle-aligned candidates are marked as such and are used
only to isolate solver reuse or rule-library usability from earlier grounding
failures.

Concretely, method-visible files are the core algorithm-input files, public
scenario-model files, and the selected rule libraries. Evaluator-only files
include source-semantic references, evaluation references, evaluator overlays,
compiled evaluator templates, and the \texttt{evaluation\_reference} portion of
paired task files.

Paired task files containing both sides are retained only as audit artifacts:
experimental runners use separated algorithm-input, scenario-model, and
evaluation-reference files, and references are loaded only after prediction.
Forbidden-key leakage checks are run on the separated files before evaluation.

The 200 tasks collectively exercise the main rule-interaction mechanisms used
by \CTHR{}, but coverage is intentionally uneven across domains and should be
read as a curated source-semantic audit rather than a balanced sampling
distribution. The architecture tasks emphasize rule applicability,
multi-rule conjunction, exception handling, precedence, and low-dimensional
regulatory thresholds. The aviation tasks include both business-realistic
multi-constraint procedure-design queries and rule-grounded cases targeting
applicability, dependency, branch exclusion, precedence, parameter propagation,
and audit traceability.

\subsubsection*{KG-to-Rule-Library Extraction Models}
Starting from the fixed KG exports, three LLM backbones are called to generate
rule libraries: qwen-plus, deepseek-v4-pro, and
mimo-v2.5-pro. Aviation and architecture use domain-specific
KG-to-rule prompt templates, which are archived with the supplementary experiment scripts. Both prompts use the same output schema: each returned rule  must contain a stable rule identifier, human-readable name, domain, rule type, source chunk IDs, optional KG node IDs, a guard, constraints, inter-rule relations, provenance fields, extraction notes, and a confidence score. The model is instructed to return only grounded JSON,
to avoid inventing numeric values or source references, and to preserve exceptions, precedence, alternatives, and dependencies as explicit rule relations.


\subsubsection*{Metrics}


CTHR and the compared methods are evaluated from five perspectives: valid-rule identification, feasible-region and decision validity, solver-backend performance, rule-library quality, and audit traceability. Unless otherwise specified, ratio-based metrics are first computed for each task and then averaged over the corresponding dataset.

\paragraph{Valid-rule identification metrics.}
To evaluate whether a method can identify the rules that are truly valid for a task, rule precision and rule recall are reported. For a task $q$, let $R_{\mathrm{pred}}(q)$ denote the predicted valid rule set and $R_{\mathrm{ref}}(q)$ denote the reference valid rule set. These metrics are defined as
\begin{equation}
\mathrm{Rule\ Precision}(q)=
\frac{|R_{\mathrm{pred}}(q)\cap R_{\mathrm{ref}}(q)|}
{|R_{\mathrm{pred}}(q)|},
\end{equation}
\begin{equation}
\mathrm{Rule\ Recall}(q)=
\frac{|R_{\mathrm{pred}}(q)\cap R_{\mathrm{ref}}(q)|}
{|R_{\mathrm{ref}}(q)|}.
\end{equation}

In the candidate-to-valid experiment, Candidate/Ref Ratio and Predicted/Ref Ratio are additionally reported. Candidate/Ref Ratio is the number of candidate rules divided by the number of reference valid rules, measuring how broad the grounding input is. Predicted/Ref Ratio is the number of predicted valid rules divided by the number of reference valid rules, measuring whether the recovered rule-set size is close to the reference. Exact Rule Match is also reported as the percentage of tasks for which the predicted rule set exactly matches the reference rule set.

\paragraph{Feasible-region and decision-validity metrics.}
To evaluate whether the optimized decisions satisfy the relevant constraint semantics, Formal CSR and Semantic CSR are reported. Formal CSR checks whether the returned decision satisfies the feasible region constructed by the evaluated method itself:
\begin{equation}
\mathrm{Formal\ CSR}=
\frac{
|\{q\in Q:\hat{x}_q\in F_{\mathrm{method}}(q)\}|
}
{|Q|}.
\end{equation}

Semantic CSR checks whether the same returned decision satisfies the hidden source-rule reference semantics:
\begin{equation}
\mathrm{Semantic\ CSR}=
\frac{
|\{q\in Q:\hat{x}_q\in F_{\mathrm{ref}}(q)\}|
}
{|Q|}.
\end{equation}

Formal CSR measures internal modeling and solver consistency, whereas Semantic CSR measures whether the returned decision is genuinely valid under the source-rule semantics. False Accept and Invalid Cases are also reported. False Accept is the percentage of tasks where the method accepts its returned decision as feasible but the reference semantics rejects it. Invalid Cases count the tasks whose final optimized decision violates the reference semantics.

\paragraph{Solver-backend metrics.}
To evaluate whether different solvers can consume the same fixed \CTHR{} feasible cells, Solve, Cell CSR, and Objective Gap are reported. Solve is the percentage of tasks for which the backend returns an optimized decision:
\begin{equation}
\mathrm{Solve}=
\frac{
|\{q\in Q:\hat{x}_q\ \mathrm{returned}\}|
}
{|Q|}.
\end{equation}

Cell CSR checks whether the returned decision satisfies the \CTHR{} feasible cell:
\begin{equation}
\mathrm{Cell\ CSR}=
\frac{
|\{q\in Q:\hat{x}_q\in P_k(q)\}|
}
{|Q|}.
\end{equation}

Objective Gap measures the relative gap between the objective value returned by a solver and the best-known objective for the same task, averaged over tasks. This metric compares optimization quality across solver backends over the same feasible cells, rather than rule-recovery quality.

\paragraph{Rule-library quality metrics.}
To compare rule libraries generated by different LLM backbones, both record-level and downstream metrics are reported. Rules is the total number of rule records generated by a given model. Provenance Valid checks whether each rule record can be traced to a valid source document, clause, or evidence chunk. In the downstream evaluation, Canonical Precision/Recall first maps model-specific rule IDs back to canonical source-rule IDs before computing rule precision and recall, avoiding false mismatches caused by different ID namespaces. Cand.-0 counts tasks for which the oracle-candidate input is empty, and Invalid counts optimized decisions that violate the source-rule semantic oracle.

\paragraph{Audit-traceability metrics.}
To evaluate auditability, Certificate Coverage, Rule Provenance Valid, and Valid-chain Trace Complete are reported. Certificate Coverage is the percentage of task-method rows for which an optimized decision is returned together with an audit certificate. Rule Provenance Valid checks whether the rule IDs listed in the certificate can be found in the rule library or KG provenance records. Valid-chain Trace Complete checks whether the certificate lists the complete valid rule structure corresponding to the active feasible cell. These metrics verify whether CTHR can trace an optimized decision back to the rule chain and original document evidence that support it.

\subsection{Overall Performance Evaluation}
\label{sec:overall_perf}

This section evaluates \CTHR{} in two settings. The first is the full
KG-to-constraint pipeline, where each method starts from candidate rules,
constructs its own feasible region, and returns an optimized decision. Flat
and native symbolic baselines consume candidate rules without the \CTHR{}
valid-rule-structure layer. The \CTHR{}-style rows use the same \CTHR{}
candidate-to-valid resolution and feasible-cell compilation, replacing only
the backend encoding and optimizer. The current 200-task evaluation reports
the numeric backends used for the new benchmark: HiGHS, CP-SAT/OR-Tools, and
SCIP. The second setting fixes the feasible cells compiled by \CTHR{} and
changes only the downstream solver, isolating backend reuse.

\begin{table*}[t]
\centering
\caption{End-to-end KG-to-constraint modeling performance on the 200-task
aviation and architecture benchmark. Each method starts from candidate rules,
constructs a feasible region, and returns an optimized decision.}
\label{tab:overall_pipeline}
\scriptsize
\setlength{\tabcolsep}{3pt}
\begin{tabularx}{\textwidth}{p{0.10\textwidth}p{0.25\textwidth}p{0.13\textwidth}cccccc}
\toprule
Dataset & Method & Type & \shortstack{Rule\\P} & \shortstack{Rule\\R} &
\shortstack{Formal\\CSR} & Sem-CSR & \shortstack{False\\accept} & Invalid \\
\midrule
Aviation & Flat baseline & Flat & 7.9\% & 91.7\% & 0.0\% & 0.0\% & 0.0\% & 150/150 \\
Aviation & Native MILP + HiGHS & Native symbolic & 2.9\% & 7.3\% & 28.0\% & 8.7\% & 19.3\% & 137/150 \\
Aviation & Native CP-SAT + OR-Tools & Native symbolic & 3.3\% & 8.3\% & 31.3\% & 10.7\% & 20.7\% & 134/150 \\
Aviation & Native SCIP & Native symbolic & 3.3\% & 8.3\% & 31.3\% & 8.7\% & 22.7\% & 137/150 \\
Aviation & \CTHR{}-style HiGHS & \CTHR{} cells & \textbf{74.7\%} & \textbf{70.7\%} & \textbf{96.7\%} & \textbf{88.7\%} & 8.0\% & \textbf{17/150} \\
Aviation & \CTHR{}-style CP-SAT + OR-Tools & \CTHR{} cells & \textbf{74.7\%} & \textbf{70.7\%} & 94.0\% & 86.0\% & 8.0\% & 21/150 \\
Aviation & \CTHR{}-style SCIP & \CTHR{} cells & \textbf{74.7\%} & \textbf{70.7\%} & \textbf{96.7\%} & \textbf{88.7\%} & 8.0\% & \textbf{17/150} \\
\midrule
Architecture & Flat baseline & Flat & 6.8\% & 99.0\% & 66.0\% & 78.0\% & 0.0\% & 11/50 \\
Architecture & Native MILP + HiGHS & Native symbolic & 51.7\% & 62.0\% & 72.0\% & 64.0\% & 8.0\% & 18/50 \\
Architecture & Native CP-SAT + OR-Tools & Native symbolic & 75.7\% & 90.0\% & \textbf{100.0\%} & 92.0\% & 8.0\% & 4/50 \\
Architecture & Native SCIP & Native symbolic & 75.7\% & 90.0\% & \textbf{100.0\%} & 92.0\% & 8.0\% & 4/50 \\
Architecture & \CTHR{}-style HiGHS & \CTHR{} cells & \textbf{95.3\%} & \textbf{95.0\%} & 96.0\% & \textbf{96.0\%} & \textbf{0.0\%} & \textbf{2/50} \\
Architecture & \CTHR{}-style CP-SAT + OR-Tools & \CTHR{} cells & \textbf{95.3\%} & \textbf{95.0\%} & 96.0\% & \textbf{96.0\%} & \textbf{0.0\%} & \textbf{2/50} \\
Architecture & \CTHR{}-style SCIP & \CTHR{} cells & \textbf{95.3\%} & \textbf{95.0\%} & 96.0\% & \textbf{96.0\%} & \textbf{0.0\%} & \textbf{2/50} \\
\bottomrule
\end{tabularx}
\end{table*}

Table~\ref{tab:overall_pipeline} shows that the flat aviation
baseline still retrieves many reference rules, with 91.7\% recall, but it
flattens incompatible and underspecified candidates into a single constraint
set and produces no source-valid decisions. Native symbolic encodings are also
weak on aviation because they must both select valid rules and encode a
heterogeneous mix of numeric, geometric, and formula-derived constraints
without the \CTHR{} valid-structure layer. Their Sem-CSR remains between
8.7\% and 10.7\%.

\CTHR{} improves the aviation result by separating rule selection from
feasible-cell compilation. The \CTHR{}-style HiGHS and SCIP rows reach 88.7\%
Sem-CSR and reduce invalid tasks from 150/150 under the flat baseline to
17/150. The CP-SAT/OR-Tools backend is slightly lower at 86.0\% Sem-CSR. These
numbers should not be read as a claim that backend choice is irrelevant:
candidate-to-valid recovery is imperfect on the 150-task aviation suite
(74.7\% precision and 70.7\% recall), and the remaining false accepts indicate
cases where the compiled cell or backend solution still fails the hidden
source-rule semantics.

The architecture benchmark presents a different pattern. Flat modeling has
very high recall (99.0\%) but very low precision (6.8\%), yielding 78.0\%
Sem-CSR and 11/50 invalid tasks. Native CP-SAT/OR-Tools and SCIP perform
strongly on this lower-dimensional building-code suite, reaching 92.0\%
Sem-CSR, but they still show 8.0\% false accept and four invalid tasks. The
\CTHR{}-style backends improve rule identification to 95.3\% precision and
95.0\% recall, remove false accepts, and reach 96.0\% Sem-CSR with 2/50 invalid
tasks. Thus, \CTHR{} is not a new optimizer; its contribution is the
source-aware modeling layer that turns KG-derived rules into feasible cells
that can then be optimized by different backends.

\begin{table*}[t]
\centering
\caption{Solver-backend evaluation over fixed \CTHR{} feasible cells. The
inputs are the latest \CTHR{} predicted valid rules and compiled constraint
templates; only the downstream solver backend is changed.}
\label{tab:cell_solver_backends}
\scriptsize
\setlength{\tabcolsep}{4pt}
\begin{tabularx}{\textwidth}{p{0.12\textwidth}Yccc}
\toprule
Dataset & Solver over \CTHR{} cells & Solve & Cell CSR & Objective gap \\
\midrule
Aviation & HiGHS over \CTHR{} cells & 96.7\% & 96.7\% & 0.0\% \\
Aviation & CP-SAT + OR-Tools over \CTHR{} cells & 96.7\% & 95.3\% & 0.7\% \\
Aviation & SCIP over \CTHR{} cells & 96.7\% & 96.7\% & 0.0\% \\
\midrule
Architecture & HiGHS over \CTHR{} cells & 100.0\% & 100.0\% & 2.0\% \\
Architecture & CP-SAT + OR-Tools over \CTHR{} cells & 100.0\% & 100.0\% & 2.0\% \\
Architecture & SCIP over \CTHR{} cells & 100.0\% & 100.0\% & 0.0\% \\
\bottomrule
\end{tabularx}
\end{table*}

Table~\ref{tab:cell_solver_backends} shows that, once \CTHR{} feasible cells
are fixed, backend behavior becomes much clearer. On architecture, all three
reported backends solve all 50 tasks and satisfy the fixed cells. On aviation,
HiGHS, CP-SAT/OR-Tools, and SCIP solve 96.7\% of the 150 tasks; HiGHS and
SCIP return only cell-satisfying solutions, while CP-SAT/OR-Tools has a
slightly lower Cell CSR of 95.3\%. Objective gaps are small for HiGHS and
SCIP on aviation, while architecture shows that SCIP obtains the tightest
objectives over the same cells. These results support the intended division of
labor: \CTHR{} defines the auditable feasible cells, and the solver backend
determines how completely and tightly those cells are optimized.

\subsection{Candidate-to-Valid Rule Recovery}
\label{sec:candidate_valid_exp}

The previous section evaluates the complete modeling-and-solving pipeline. This
section isolates the rule-recovery step before any feasible-region compilation
or downstream optimization is considered. The goal is to test whether \CTHR{}
can recover the valid rules for a grounded scenario from the candidate set. The
broad candidate set is constructed using only method-visible task fields,
public scenario facts, decision variables, rule metadata, guards, relation
fields, provenance, and source-domain hints. Reference valid rules, hidden
constraints, reference feasible cells, semantic validators, and solver outcomes
are not exposed to candidate grounding, profile filtering, the optional LLM
reranker, or the final relation-selection step; they are used only after
prediction for evaluation. In this experiment, \CTHR{} uses the current
candidate-to-valid profile resolver. The aviation run uses the aviation recall
guard without online LLM filtering, while the architecture run uses the cached
profile-auto resolver.

A set of candidate-to-valid rule-identification metrics is used to evaluate whether CTHR can identify task-valid rules from candidate rules. Candidate/Ref is the size of the input broad candidate set $C(q)$ in Algorithm~\ref{alg:rule_interaction_resolution}, normalized by the number of reference valid rules. Filtered/Ref is the size of the profile-constrained set $C_{\mathrm{prof}}(q)$ obtained after lines 1--14, normalized by the same reference count. Predicted/Ref is the size of the final valid rule set $V(q)$ returned by Algorithm~\ref{alg:rule_interaction_resolution}, again normalized by the reference count. Rule-ID precision and recall measure the set overlap between predicted and reference rule identifiers. Exact match is the percentage of tasks for which the predicted and reference rule sets are identical. Extra and Missing count non-reference rules included in the prediction and reference rules missed by the method, respectively.

\begin{table*}[t]
\centering
\caption{Aggregate candidate-to-valid rule recovery results. Candidate/Ref
measures the breadth of the method-visible grounded candidate set, Filtered/Ref
measures the intermediate profile-constrained candidate set, and Predicted/Ref
measures the final \CTHR{}-selected valid-rule set, all normalized by the
reference valid-rule count used only for evaluation.}
\label{tab:candidate_valid_aggregate}
\scriptsize
\setlength{\tabcolsep}{3pt}
\begin{tabularx}{\textwidth}{Yrrrrrrrrr}
\toprule
Dataset & Tasks & \shortstack{Candidate\\ / Ref} &
\shortstack{Filtered\\ / Ref} & \shortstack{Predicted\\ / Ref} &
\shortstack{Rule-ID\\ precision} & \shortstack{Rule-ID\\ recall} &
\shortstack{Exact\\ match}  \\
\midrule
Aviation & 150 & 11.237 & 1.413 & 0.960 & 0.747 & 0.707 & 0.700 \\
Architecture & 50 & 15.387 & 1.670 & 1.340 & 0.933 & 0.930 & 0.900  \\
\bottomrule
\end{tabularx}
\end{table*}

Table~\ref{tab:candidate_valid_aggregate} shows that the candidate sets are
substantially broader than the reference valid-rule sets. Aviation starts from
11.237 times as many candidate rules as reference rules on average. \CTHR{}
reduces the intermediate Filtered/Ref ratio to 1.413 and the final
Predicted/Ref ratio to 0.960, but the enlarged aviation benchmark remains
challenging: Rule-ID precision is 74.7\%, Rule-ID recall is 70.7\%, and the
latest exact-match audit is 70.0\%.

Architecture begins from an even wider candidate pool, 15.387 times the
reference size, but the profile-auto resolver contracts it to a Filtered/Ref
ratio of 1.670 and a final Predicted/Ref ratio of 1.340. It maintains 93.3\%
Rule-ID precision, 93.0\% Rule-ID recall, and 90.0\% exact match. The architecture audit also reports 70.0\% semantic constraint satisfaction and 92.0\% formal feasibility for this
isolated candidate-to-valid run, indicating that some failures after rule
identification arise from downstream solving or semantic checking rather than
from rule-ID recovery alone.
\FloatBarrier

\subsection{Rule-Library Generator Comparison}
\label{sec:rule_library_exp}

The previous experiments use the qwen-plus generated rule library as the
default rule library for each domain. This section therefore treats qwen-plus as
the main rule-library source and uses DeepSeek-v4-Pro and Xiaomi MIMO-V2.5-Pro as two
additional replacement libraries to test whether the modeling layer can consume
rules produced by different LLMs. The evaluation has two levels. First, the
generated rule records themselves are audited in terms of rule count,
provenance validity, constraint grounding, and relation grounding. Second, each
model-specific rule library is substituted into a downstream \CTHR{} run while keeping task
parameters, source-rule semantic oracle, and solver backend fixed. This
experiment is diagnostic: it tests rule-library usability under controlled
alignment and compilation support, not the success rate of a fully automatic
KG-to-rule extraction pipeline.

\begin{table*}[t]
\centering
\caption{Quality audit of LLM-generated rule libraries. The Rules column counts
all rule records emitted by the generator. Provenance validity checks whether
each rule can be traced to the expected source document, clause, or evidence
chunk.}
\label{tab:rule_library_quality}
\footnotesize
\begin{tabular*}{\textwidth}{@{}l@{\extracolsep{\fill}}lrrrr@{}}
\toprule
Dataset & Generator & Rules & \shortstack{Provenance\\valid} & \shortstack{Constraint\\grounding} & \shortstack{Relation\\grounding} \\
\midrule
Aviation & qwen-plus & 111 & 100.0\% & 100.0\% & 100.0\% \\
Aviation & DeepSeek-v4-Pro & 142 & 100.0\% & 99.1\% & 99.6\% \\
Aviation & MIMO-V2.5-Pro & 186 & 100.0\% & 100.0\% & 100.0\% \\
\midrule
Architecture & qwen-plus & 1,585 & 100.0\% & 100.0\% & 100.0\% \\
Architecture & DeepSeek-v4-Pro & 2,419 & 100.0\% & 99.5\% & 99.6\% \\
Architecture & MIMO-V2.5-Pro & 2,157 & 99.9\% & 99.6\% & 99.8\% \\
\bottomrule
\end{tabular*}
\end{table*}

Table~\ref{tab:rule_library_quality} shows that the three generators all
produce largely well-grounded rule records. Provenance validity is at least
99.9\% across all six domain-generator pairs, and constraint and relation
grounding are also above 99.1\%. The generators differ more in rule-library
size and downstream usability than in whether the emitted records contain
traceable source references.
\FloatBarrier

\begin{table*}[t]
\centering
\caption{Downstream performance when replacing the rule library generator under
a controlled \CTHR{} setting. The solver backend is fixed to SCIP, matching the
\CTHR{}-style SCIP rows in Table~\ref{tab:overall_pipeline}. Rule P/R, Formal
CSR, Sem-CSR, and False accept are reported as percentages. Cand.-0 counts tasks
with no candidate rules, and Valid-0 counts tasks for which no valid rules
remain after \CTHR{} recovery.}
\label{tab:rule_library_downstream}
\scriptsize
\setlength{\tabcolsep}{2.4pt}
\renewcommand{\arraystretch}{0.92}
\begin{tabular*}{\textwidth}{@{}l@{\extracolsep{\fill}}lclcccccc@{}}
\toprule
Dataset & Generator & Backend & Rule P/R & Formal CSR &
Sem-CSR & False acc. & Cand.-0 & Valid-0 & Invalid \\
\midrule
Aviation & qwen-plus & SCIP & 74.7\%/70.7\% & 96.7\% & 88.7\% & 8.0\% & 0 & 0 & 17/150 \\
Aviation & DeepSeek-v4-Pro & SCIP & 76.8\%/70.0\% & 96.7\% & 86.7\% & 0.0\% & 0 & 0 & 15/150 \\
Aviation & MIMO-V2.5-Pro & SCIP & 78.9\%/70.0\% & 94.0\% & 84.0\% & 0.0\% & 0 & 0 & 19/150 \\
\midrule
Architecture & qwen-plus & SCIP & 95.3\%/95.0\% & 96.0\% & 96.0\% & 0.0\% & 0 & 0 & 2/50 \\
Architecture & DeepSeek-v4-Pro & SCIP & 96.0\%/100.0\% & 96.0\% & 96.0\% & 0.0\% & 0 & 0 & 2/50 \\
Architecture & MIMO-V2.5-Pro & SCIP & 92.0\%/100.0\% & 96.0\% & 96.0\% & 0.0\% & 0 & 0 & 2/50 \\
\bottomrule
\end{tabular*}
\end{table*}

Because different LLMs may express the same source rule with different
granularities, rule names, and textual descriptions, direct evaluation with
model-specific rule IDs can introduce alignment failures. The evaluation therefore maps
model-specific rules to canonical source-rule identifiers before computing Rule
P/R. Table~\ref{tab:rule_library_downstream} uses the same SCIP backend for all
rows. The qwen-plus rows are the default rule-library results and therefore
match the corresponding \CTHR{}-style SCIP rows in
Table~\ref{tab:overall_pipeline}. DeepSeek-Pro and Xiaomi MIMO are replacement
libraries evaluated under the same backend and task set. In aviation, all three
libraries avoid candidate-empty and valid-rule-empty failures. The default qwen-plus 
reaches 88.7\% Sem-CSR with 17/150 invalid tasks. DeepSeek-Pro achieves a higher 
rule precision of 76.8\% and a comparable recall of 70.0\%, yielding 86.7\% Sem-CSR 
and 15/150 invalid cases. Xiaomi MIMO exhibits the highest rule precision at 78.9\%, 
but its Sem-CSR drops to 84.0\% with 19/150 invalid cases. In architecture, all 
three libraries similarly avoid candidate-empty and valid-rule-empty failures, 
uniformly reaching 96.0\% Sem-CSR with exactly 2/50 invalid tasks, varying only 
slightly in their rule precision and recall granularities. These results indicate 
that \CTHR{} can consume multiple generated rule libraries under controlled alignment, 
but downstream validity still depends on generator-specific rule granularity 
and executable semantics.
\FloatBarrier


\subsection{Audit Traceability}
\label{sec:audit_traceability_exp}

The preceding experiments evaluate whether the returned decisions are feasible
and source-consistent. This section evaluates a different requirement: whether a
decision can be audited after optimization. General-purpose solvers such as
HiGHS, CP-SAT/OR-Tools, and SCIP can return feasible assignments, but they
do not by themselves explain which source-linked rule chain was selected for
the accepted decision. In \CTHR{}, the solver is used only after the valid rule
structures have been compiled into traceable feasible cells. Therefore, the
audit certificate is a \CTHR{} metadata object generated from the selected
feasible cell and its rule chain, rather than from a solver-specific proof log.

Three backend replacements are evaluated over the same \CTHR{}
feasible cells: HiGHS, CP-SAT/OR-Tools, and SCIP. The evaluation is conducted
on the new 200-task benchmark. \emph{Certificate coverage} is defined as the
fraction of task-backend rows for which a returned cell-level solution is
accompanied by an audit certificate. \emph{Rule provenance validity} checks
whether each rule identifier listed in the certificate resolves to a rule
record in the current rule library and whether that record contains non-empty
provenance metadata such as document, section, or KG chunk identifiers.
\emph{Valid-chain trace completeness} checks whether the certificate lists all
rule identifiers in the selected \CTHR{} valid-rule chain used to compile the
active feasible cell. Consequently, an accepted solution can be traced from the
audit certificate to the valid rule chain, the rule-library entries, and the
document or KG metadata stored in their provenance fields.

\begin{table*}[t]
\centering
\caption{Audit traceability of optimized decisions returned through different solver backends on the 200-task benchmark. Certificate coverage and valid-chain trace completeness are measured over task-backend rows; rule provenance validity is measured over the rule identifiers appearing in generated certificates.}
\label{tab:audit_traceability}
\footnotesize
\begin{tabular*}{\textwidth}{@{}l@{\extracolsep{\fill}}lcccc@{}}
\toprule
Dataset & Method & \shortstack{Certificate\\coverage} & \shortstack{Avg. rules\\per cert.} & \shortstack{Rule provenance\\valid} & \shortstack{Valid-chain trace\\complete} \\
\midrule
Aviation & \CTHR{}+HiGHS & 84.7\% & 1.906 & 100.0\% & 84.7\% \\
Aviation & \CTHR{}+CP-SAT/OR-Tools & 95.3\% & 1.916 & 100.0\% & 95.3\% \\
Aviation & \CTHR{}+SCIP & 96.7\% & 1.917 & 100.0\% & 96.7\% \\
Architecture & \CTHR{}+HiGHS & 60.0\% & 2.233 & 100.0\% & 60.0\% \\
Architecture & \CTHR{}+CP-SAT/OR-Tools & 88.0\% & 2.045 & 100.0\% & 88.0\% \\
Architecture & \CTHR{}+SCIP & 88.0\% & 2.045 & 100.0\% & 88.0\% \\
\bottomrule
\end{tabular*}
\end{table*}

Table~\ref{tab:audit_traceability} shows that, whenever a backend returns a
valid \CTHR{} cell-level solution, the generated certificate preserves the
rule-chain and provenance links. All certificate rule identifiers in both
domains resolve to provenance-bearing rule-library records. The main difference
across backends is coverage: HiGHS covers 84.7\% of aviation tasks but only
60.0\% of architecture tasks, while CP-SAT/OR-Tools and SCIP cover 95.3\% and
96.7\% of aviation tasks and 88.0\% of architecture tasks. This pattern is
consistent with the cell-solver results: missing certificates mainly arise
when the backend does not return a valid active-cell solution under the current
numeric feasibility check, rather than from a broken provenance lookup. These
results should be interpreted as a machine-checkable rule-chain and provenance
linkage test, not as a solver-native proof log or a human audit-time study.
They support the intended division of labor in \CTHR{}: external solvers
optimize over feasible cells, while the \CTHR{} compilation layer supplies the
auditable link from an optimized decision back to the KG-derived rules and
source evidence.
\FloatBarrier



\subsection{Component Ablation}
\label{sec:ablation_exp}

This experiment ablates key components of the \CTHR{} modeling layer on the 200-task benchmark. The default configuration uses the predicted valid
rules produced by the candidate-to-valid stage, the compiled
rule-to-constraint templates, and the same backend family used in the main
experiments. Each variant removes one component: candidate-to-valid profiling,
scenario applicability, positive closure, negative resolution,
relation-aware recovery, or compiled templates. The experiment reports solving success,
formal and semantic constraint satisfaction, Rule-ID precision/recall. This is a component-ablation setting: variants remove
selected modeling components rather than using hidden labels as inputs.

\begin{table*}[t]
\centering
\caption{Component ablation of the \CTHR{} pipeline on the new aviation and architecture benchmarks. Rule P/R denotes per-task Rule-ID precision and recall; Exact denotes exact rule-set match. Aviation uses the SCIP backend version aligned with the main end-to-end evaluation.}
\label{tab:component_ablation}
\scriptsize
\begin{tabular*}{\textwidth}{@{}l@{\extracolsep{\fill}}lccccc@{}}
\toprule
Dataset & Variant & Solved & \shortstack{Formal\\CSR} & Sem-CSR & \shortstack{Rule\\P/R}  \\
\midrule
Aviation & \CTHR{} default + SCIP & 133/150 & 96.7\% & 88.7\% & 74.7\% / 70.7\%  \\
Aviation & w/o Candidate-to-Valid Profiling & 34/150 & 22.7\% & 22.7\% & 8.8\% / 88.7\%  \\
Aviation & w/o Scenario Applicability & 31/150 & 20.7\% & 20.7\% & 8.8\% / 88.7\%  \\
Aviation & w/o Positive Closure & 34/150 & 22.7\% & 22.7\% & 8.8\% / 88.7\%  \\
Aviation & w/o Negative Resolution & 34/150 & 22.7\% & 22.7\% & 8.9\% / 91.7\%  \\
Aviation & w/o Relation-Aware Recovery & 128/150 & 93.3\% & 85.3\% & 62.7\% / 81.0\%  \\
Aviation & w/o Compiled Templates & 113/150 & 83.3\% & 0.0\% & 74.7\% / 70.7\%  \\
\midrule
Architecture & \CTHR{} default + SCIP & 48/50 & 86.0\% & 72.0\% & 81.5\% / 95.0\%  \\
Architecture & w/o Candidate-to-Valid Profiling & 48/50 & 86.0\% & 72.0\% & 81.5\% / 95.0\%  \\
Architecture & w/o Scenario Applicability & 50/50 & 90.0\% & 76.0\% & 77.5\% / 97.0\% \\
Architecture & w/o Positive Closure & 48/50 & 86.0\% & 72.0\% & 81.5\% / 95.0\%  \\
Architecture & w/o Negative Resolution & 48/50 & 86.0\% & 72.0\% & 80.1\% / 95.0\%  \\
Architecture & w/o Relation-Aware Recovery & 50/50 & 90.0\% & 76.0\% & 77.1\% / 99.0\%  \\
Architecture & w/o Compiled Templates & 48/50 & 80.0\% & 66.0\% & 81.5\% / 95.0\%  \\
\bottomrule
\end{tabular*}
\end{table*}
\FloatBarrier

\begin{table*}[t]
\centering
\caption{Relation-level ablation of the \CTHR{} rule-interaction layer. Each row removes one relation type from relation-aware recovery while keeping the SCIP backend fixed. Rule P/R denotes per-task Rule-ID precision and recall; Exact denotes exact rule-set match.}
\label{tab:relation_ablation}
\scriptsize
\setlength{\tabcolsep}{3.5pt}
\begin{tabular*}{\textwidth}{@{}l@{\extracolsep{\fill}}lcccc@{}}
\toprule
Dataset & Removed relation & Solved & Sem-CSR & Rule P/R & Exact \\
\midrule
Aviation & None & 133/150 & 88.7\% & 74.7\% / 70.7\% & 132/150 \\
Aviation & Scenario applicability & 31/150 & 20.7\% & 8.8\% / 88.7\% & 0/150 \\
Aviation & Dependency & 128/150 & 85.3\% & 62.7\% / 81.0\% & 122/150 \\
Aviation & Exclusion / branch & 128/150 & 85.3\% & 62.7\% / 81.0\% & 122/150 \\
Aviation & Override & 128/150 & 85.3\% & 62.7\% / 81.0\% & 122/150 \\
Aviation & Precedence & 128/150 & 85.3\% & 62.7\% / 81.0\% & 122/150 \\
Aviation & Parameter propagation & 128/150 & 85.3\% & 62.7\% / 81.0\% & 122/150 \\
\midrule
Architecture & None & 48/50 & 72.0\% & 81.5\% / 95.0\% & 37/50 \\
Architecture & Scenario applicability & 50/50 & 76.0\% & 77.5\% / 97.0\% & 33/50 \\
Architecture & Dependency & 48/50 & 72.0\% & 81.5\% / 95.0\% & 37/50 \\
Architecture & Exclusion / branch & 48/50 & 72.0\% & 80.1\% / 95.0\% & 35/50 \\
Architecture & Override & 48/50 & 72.0\% & 81.4\% / 95.0\% & 37/50 \\
Architecture & Precedence & 48/50 & 72.0\% & 81.5\% / 95.0\% & 37/50 \\
Architecture & Parameter propagation & 48/50 & 72.0\% & 81.5\% / 95.0\% & 37/50 \\
\bottomrule
\end{tabular*}
\end{table*}
\FloatBarrier



Table~\ref{tab:component_ablation} shows that the component-level effects differ
substantially across the two domains. In aviation, removing the early
candidate-to-valid profile, scenario applicability, or closure/resolution
components collapses the end-to-end solve rate to 20.7\%--22.7\% and reduces
Rule-ID precision to below 9\% (8.8\%--8.9\%), even though recall remains high 
(88.7\%--91.7\%). This confirms that broad aviation grounding contains many rules 
that are related to the procedure-design corpus but should not jointly define 
the current feasible cell. Removing relation-aware recovery is less destructive 
but drops the solve rate to 128/150, reducing Rule-ID precision to 62.7\% while 
increasing recall to 81.0\%. Removing compiled templates keeps the same 
valid-rule identifiers but drives Sem-CSR to 0.0\%, isolating the 
constraint-compilation layer as essential for executable source semantics. The
architecture ablation is less severe because the new architecture tasks and
profile resolver already produce cleaner candidate sets, but removing scenario
applicability or relation-aware recovery still increases recall up to 97.0\%--99.0\% 
at the cost of lower precision, indicating a looser rule set.

Table~\ref{tab:relation_ablation} reports the companion relation-level
ablation. In aviation, removing scenario applicability is the strongest
relation-level failure mode: the solve rate falls to 31/150, Sem-CSR to 20.7\%,
Rule P/R to 8.8\%/88.7\%, and exact match to 0/150. Removing dependency,
exclusion/branch, override, precedence, or parameter propagation gives an
identical degradation pattern, with 128/150 solved tasks, 85.3\% Sem-CSR,
62.7\%/81.0\% Rule P/R, and 122/150 exact matches. In architecture, removing
scenario applicability and exclusion/branch mainly affects rule-set compactness:
exact match drops from 37/50 to 33/50 and 35/50, respectively, whereas
dependency, precedence, and parameter propagation remain marginal on this 50-task
subset, preserving perfect alignment with the default baseline. Overall, the 
ablation supports a layered necessity claim: early profiling and scenario 
applicability control candidate noise, relation-aware recovery improves 
rule-set compactness, and compiled templates are needed to preserve 
source semantics at the constraint level.

\FloatBarrier

\subsection{Fairness and Controlled Scalability Checks}
\label{sec:fairness_scalability_exp}

This section reports two controlled supplemental checks. First, all
model-rule libraries are evaluated on the same task-level membership probes,
including midpoint, boundary, near-boundary, random, and reference-feasible
anchor points, with labels computed from hidden source-semantic references.
The probe count is the total number of shared membership-test points, not the
number of tasks or optimized solutions. Second, irrelevant candidate rules are
injected to expand the candidate set to 2, 4, and 8 times its original size.
These checks evaluate probe-design fairness and controlled candidate-noise
robustness; they are not industrial-scale deployment claims.

\begin{table*}[t]
\centering
\caption{Fairness and controlled candidate-noise checks. Panel (a) evaluates model-specific rule libraries on shared membership probes; panel (b) injects irrelevant candidate rules and reports aggregated robustness.}
\label{tab:fairness_scalability}
\begingroup
\scriptsize
\setlength{\tabcolsep}{2.5pt}
\renewcommand{\arraystretch}{0.9}
\textbf{(a) Unified membership-probe evaluation}\\[-2pt]
\begin{tabular*}{\textwidth}{@{}l@{\extracolsep{\fill}}lccccc@{}}
\toprule
Dataset & Model-rule library & \shortstack{Probe\\points} &
\shortstack{Semantic\\consistency} & \shortstack{False\\accept} &
\shortstack{False\\reject} & \shortstack{Accepted-point\\Sem-CSR} \\
\midrule
Aviation & DeepSeek-v4-Pro & 25,253 & 99.8\% & 0.0\% & 0.2\% & 98.5\% \\
Aviation & qwen-plus & 25,253 & 100.0\% & 0.0\% & 0.0\% & 100.0\% \\
Aviation & Xiaomi MIMO & 25,253 & 99.6\% & 0.0\% & 0.4\% & 100.0\% \\
Architecture & DeepSeek-v4-Pro & 1,766 & 100.0\% & 0.0\% & 0.0\% & 100.0\% \\
Architecture & qwen-plus & 1,766 & 99.9\% & 0.1\% & 0.0\% & 99.4\% \\
Architecture & Xiaomi MIMO & 1,766 & 99.3\% & 0.7\% & 0.0\% & 96.4\% \\
\bottomrule
\end{tabular*}
\vspace{3pt}
\textbf{(b) Controlled candidate-noise stress test}\\[-2pt]
\begin{tabular*}{\textwidth}{@{}l@{\extracolsep{\fill}}lcccccc@{}}
\toprule
Dataset & \shortstack{Candidate\\factor} &
\shortstack{Mean\\candidate rules} & \shortstack{Mean compiled\\constraints} &
\shortstack{Mean semantic\\consistency} & \shortstack{Minimum semantic\\consistency} &
\shortstack{Maximum false\\accept} & \shortstack{Maximum false\\reject} \\
\midrule
Aviation & 1$\times$ & 2.173 & 48.649 & 99.8\% & 99.5\% & 0.0\% & 0.5\% \\
Aviation & 2$\times$ & 4.347 & 49.056 & 99.6\% & 99.4\% & 0.0\% & 0.6\% \\
Aviation & 4$\times$ & 8.693 & 50.031 & 99.3\% & 99.2\% & 0.0\% & 0.8\% \\
Aviation & 8$\times$ & 17.387 & 52.151 & 98.9\% & 98.6\% & 0.0\% & 1.4\% \\
Architecture & 1$\times$ & 1.907 & 11.333 & 100.0\% & 100.0\% & 0.0\% & 0.0\% \\
Architecture & 2$\times$ & 3.813 & 11.493 & 100.0\% & 100.0\% & 0.0\% & 0.0\% \\
Architecture & 4$\times$ & 7.627 & 11.593 & 99.9\% & 99.7\% & 0.0\% & 0.3\% \\
Architecture & 8$\times$ & 15.253 & 11.953 & 100.0\% & 100.0\% & 0.0\% & 0.0\% \\
\bottomrule
\end{tabular*}
\endgroup
\end{table*}

Table~\ref{tab:fairness_scalability} shows that the shared-probe evaluation
does not favor a particular model-rule library. Semantic consistency ranges
from 99.6\% to 100.0\% in aviation and from 99.3\% to 100.0\% in architecture;
the largest false-accept rate is 0.7\%. Under controlled candidate-noise
injection, aviation remains above 98.9\% mean semantic consistency at the
8$\times$ scale, and architecture remains between 99.9\% and 100.0\%. The
largest observed false-reject rate is 1.4\%.

\FloatBarrier


%=============================================================================
\section{Discussion and Limitations}
\label{sec:discussion_limitations}

\CTHR{} should be interpreted as a provenance-preserving modeling layer, not as
a replacement for optimization solvers. Its contribution is to construct
source-consistent feasible cells and audit certificates before a backend solver
searches for an objective-optimal point. The backend still matters for
optimization strength, as reflected by the objective-gap results in
Table~\ref{tab:cell_solver_backends}. The solver-backend experiments therefore
demonstrate representative adapters rather than universal solver coverage:
additional SAT, SMT, MILP, NLP, global-optimization, or heuristic solvers would
require their own encodings of the emitted feasible cells and objectives.

Some benchmark instances still do not produce a certified valid optimized
decision. This limitation should be distinguished from the semantic preservation
claim: the theory is conditional on correct valid-rule recovery and correct
rule-to-constraint compilation, whereas the end-to-end experiments also include
candidate filtering errors, rule-library imperfections, and backend-specific
modeling limits. In the aviation benchmark, the remaining failures are mostly
observed as source-semantically invalid returned decisions or missing backend
solutions on a subset of tasks, despite the large reduction in invalid cases
relative to flat and native symbolic baselines. In the architecture benchmark,
the number of such failures is smaller, but several tasks still fail because
selected rule chains can be empty, overly broad, or not fully supported by a
particular backend encoding. These cases indicate that \CTHR{} supplies an
auditable feasible-cell modeling interface for returned solutions, but it does
not by itself guarantee that every source-grounded query is solvable by every
backend. Improving fallback formulations, backend-specific encodings for
indicator or count predicates, and calibration of candidate-to-valid recovery
remain important future work.

The benchmark is a systematic two-domain task suite rather than a full
industrial deployment. It supports controlled evaluation of source-grounded rule
recovery, feasible-cell construction, decision validity, and certificate
traceability, but broader deployment would require larger corpora and further
expert review. The shared-probe and candidate-noise checks are controlled
stress tests on this benchmark, not evidence of industrial-scale deployment.
The clean split remains curated: synthetic aviation stress rules
are excluded, architecture has limited multi-cell defeasible-chain stress, and
some building-code requirements are encoded as executable semantic indicator or
count predicates. Rule-library construction depends on KG-to-rule generator
quality; downstream usability varies across LLM backbones and relies on exact
or strong rule-ID alignment overlays. Finally, the architecture benchmark
retains a small number of invalid cases, and certificate evaluation checks
machine-readable traceability rather than human audit time. The current version also does not
report an independent blind annotation study or inter-annotator agreement for
the source-semantic references in this version.

%=============================================================================
\section{Conclusion}
\label{sec:conclusion}

This paper presents \CTHR{}, a provenance-preserving KG-to-rule-to-feasible-cell
modeling layer for auditable KG-driven optimization. \CTHR{} resolves
defeasible rule interactions, constructs valid rule structures, compiles them
into source-semantics-preserving feasible cells, and returns audit certificates
that link optimized decisions back to rule and document provenance.

Across the 200-task aviation and architecture benchmark, \CTHR{} improves
source-semantic decision validity over flat-rule modeling and supports multiple
solver backends over the same feasible cells. The results show that auditable
optimization depends not only on solver strength, but also on the upstream
modeling step that decides which source rules jointly define the feasible
region. Future work will extend KG-to-rule extraction, broaden the benchmark to
additional regulated engineering domains, and reduce manual curation when
natural-language design requests are converted into structured scenarios.

%=============================================================================
\section*{Funding}
This research was supported by the National Natural Science Foundation of China
(NSFC) under Grant No.~625B2123 and the Pioneer Fund of the System of Systems
and Artificial Intelligence Key Laboratory of Sichuan Province under Grant
No.~SSAI202506.

%=============================================================================
\appendix
\setcounter{table}{0}
\renewcommand{\thetable}{A.\arabic{table}}

\section{Aviation Task Expert Review}
\label{app:aviation_expert_review}

To audit the quality of the aviation benchmark and its hidden reference labels,
100 aviation tasks were sampled from the 150-task aviation suite for expert
review. The review was conducted by three flight-procedure design experts who
had passed the CAAC Flight Procedure Design Basic Knowledge Examination. The
audit was performed after benchmark construction and was used only to assess
task and reference-label quality; it was not provided to \CTHR{}, the candidate
grounder, the profile resolver, or any solver backend.

Each expert scored every sampled task on five criteria using a 1--5 scale:
5 means that the task can be retained directly, 4 means that it is basically
reasonable with only minor issues, 3 means that it is usable after minor
revision, 2 means that it has clear issues and requires restructuring, and
1 means that it is unsuitable and should be removed or replaced. Table~\ref{tab:aviation_expert_review}
reports the mean scores across experts and reviewed tasks.

\begin{table}[H]
\centering
\caption{Expert-review protocol for the 100-task aviation benchmark sample. Scores are averaged over three certified flight-procedure design experts and the reviewed tasks.}
\label{tab:aviation_expert_review}
\footnotesize
\setlength{\tabcolsep}{2pt}
\begin{tabularx}{\columnwidth}{p{0.08\columnwidth}Xc}
\toprule
ID & Evaluation question and review focus & Mean \\
\midrule
Q1 & Scenario realism: whether the scenario resembles a realistic flight-procedure design, regulation-review, or engineering-coordination task. & 4.5 \\
Q2 & Reference-rule rationality: whether valid rules, excluded rules, priority or exception relations, and feasible-region constraints are reasonable. & 4.0 \\
Q3 & Parameter and variable plausibility: whether units, ranges, variable names, and numerical magnitudes match domain practice. & 4.2 \\
Q4 & Objective trade-off realism: whether the objectives reflect a meaningful operational trade-off rather than redundant goals. & 4.3 \\
Q5 & Leakage avoidance: whether public inputs avoid exposing rule identifiers, reference answers, feasible-region labels, or certificate targets. & 4.7 \\
\bottomrule
\end{tabularx}
\end{table}

\section{Proofs}
\label{app:proofs}

\subsection{Proof of Lemma~L1 (Applicability Correctness)}

\begin{proof}
The compiler constructs $\gamma_k$ as the conjunction of:
(a) grounded antecedents $\phi_r(\omega)$ for all $r\in\chain(k)$;
(b) dependency closures: for each $r\in\chain(k)$, all $r'\in\textsf{dep}^+(r)$
    (transitive dependency closure) must also be applicable;
(c) internal exclusion consistency: no pair $r,r'\in\chain(k)$ may satisfy
    $(r,r')\in\textsf{excl}$.
Unsatisfiable guards are removed during Z3 normalization.
This condition does not require an excluded alternative outside $\chain(k)$ to
be inapplicable. If two mutually exclusive alternatives are both applicable,
they form separate source-level extensions and therefore separate feasible
cells. The remaining guards are extensionally equivalent to source-level
applicability by construction of items (a)--(c).  
\end{proof}

\subsection{Proof of Lemma~L2 (Override Correctness)}

\begin{proof}
When a conflict class is resolved, the compiler identifies all applicable rules
in the class and applies the defeat relation.
For each $r_e \xrightarrow{\text{can\_override}} r_b$ with both applicable:
the compiler marks $r_b$ as defeated and removes its consequent before polytope emission.
Hence the emitted constraints for cell $k$ contain $c_{r_e}$ and exclude $c_{r_b}$,
matching override semantics rather than flat conjunction.  
\end{proof}

\subsection{Proof of Lemma~L3 (Precedence Correctness)}

\begin{proof}
Cross-document specificity, category specificity, and explicit superiority relations
are compiled into a precedence order $\succ_\omega$ over conflicting rules.
The compiler selects maximal non-defeated valid rule structures under this order using topological
processing of the precedence DAG (well-founded by Assumption~\ref{ass:defeat}).
Non-conflicting applicable rules are in disjoint conflict classes and are both retained.
Therefore chain selection is extensionally equivalent to source-level precedence resolution.  
\end{proof}

\subsection{Proof of Lemma~L4 (Parameter Propagation Correctness)}

\begin{proof}
By Assumption~\ref{ass:dep}, the \textsf{dep} relation is acyclic, so $\eta(\omega)$
is uniquely determined by topological evaluation.
The compiler performs symbolic substitution of $\eta(\omega)$ into
each rule consequent $c_r(x,\eta) = \langle a_r(\eta),x\rangle \leq b_r(\eta)$
before affine canonicalization.
Since $\eta(\omega)$ is unique and substitution is exact (no rounding),
$c_r(x,\eta(\omega)) \iff \hat{c}_{r,\omega}(x)$ pointwise.  
\end{proof}

\subsection{Proof of Lemma~L5 (Geometry Correctness)}

\begin{proof}
$P_k$ is defined as the H-representation produced by stacking exactly the
surviving instantiated inequalities of $\chain(k)$ after defeasible resolution.
Canonical affine rewrites preserve constraint meaning.
Hence $x\in P_k \iff x$ satisfies all surviving instantiated consequents of $\chain(k)$.  
\end{proof}

%=============================================================================
\bibliographystyle{elsarticle-num}
\bibliography{cthr_refs}

\end{document}
